\chapter{One-loop components along the supercharges, the special supercharges and the special conformal generators: design}

The components of the BV differential of an observable along the odd generators of $\mathfrak{psl}(4|4)$ and along the special conformal fields carry the one-loop corrections that the user calls $Q_1$, $S_1$ and $K_1$; along the dilatation they carry $D_1$. These generators move the regulator [ch07, prop:c7-dil, prop:c7-susy], so their one-loop parts are finite local functionals produced at the UV boundary. The observables are built at the endpoint $T$ of the transfer of [C07], with the propagator $P(\varepsilon,T)$; the components are local representatives, the same at every $T$, and $T$ is sent to $\infty$ wherever this is safe. Two questions of principle arise: the holomorphic-twist $Q_1$ of the literature must come out of the untwisted cochain without a twist in the input, and the special conformal field grows quadratically at large distance, so that its internal insertions sit at the borderline $|y|^{-4}$ of integrability [ch07, rem:c7-ir] and produce logarithms of $T$. This chapter defines the components at the transfer endpoint, lists the graphs, derives the twisted $Q_1$ as the evaluation of the untwisted cochain on a rank-one supercharge, writes the transformation laws under every scheme choice, including the outer derivation of $\mathfrak{psl}(4|4)$, together with the scheme-independent content, sets up the large-distance analysis and the definition of $K_1$ at finite $T$, and fixes the order of the engine computations.

\section{Conventions and names}

\begin{definition}[Spinor parameters, generators, names]
\label{def:o-names}
\begin{itemize}
\item Spinor parameters: $\xi\in\Sigma$ with two-component form $(\xi_\alpha{}^i,\bar\xi_{\dot\alpha i})$, $\eta\in\Sigma^*$ with two-component form $(\eta^\alpha{}_i,\bar\eta^{\dot\alpha i})$ [dictionary, Thm thm:real]. They are the parameters called $Q_1\in\Sigma$, $Q_2\in\Sigma^*$ in [ch02, def:c2-odd]; those two symbols do not occur in this chapter.
\item The holomorphic supercharge: $\xi=u\otimes w\in S_+\otimes W$ with $u\in S_+$, $w\in W$ nonzero; two-component $\xi_\alpha{}^i=u_\alpha w^i$, $\bar\xi_{\dot\alpha i}=0$. Fixed auxiliary data: $u'\in S_+$ with $u'^\alpha u_\alpha=1$ and $n\in W^*$ with $n_iw^i=1$ [targets, §2.1]. The letter $w$ without subscript always denotes this element of $W$; $w_1,w_2$ with subscripts are elements of $W_{\mathbb R}^\perp$ as in [ch05, def:c5-bilinear].
\item The conformal Killing spinor of $\xi+\eta$: $x\mapsto\xi+\rho(x)\eta$ [ch02, def:c2-odd]; two-component (K1) of the dictionary:
\[
\epsilon_\alpha{}^i(x)=\xi_\alpha{}^i-i\,x_m\sigma^m_{\alpha\dot\alpha}\bar\eta^{\dot\alpha i},\qquad \bar\epsilon_{\dot\alpha i}(x)=\bar\xi_{\dot\alpha i}+i\,x_m\eta^\alpha{}_i\sigma^m_{\alpha\dot\alpha}.
\]
\item $K_{v_1}$: the special conformal field $x\mapsto(x,x)v_1-2(v_1,x)x$, $v_1\in V$ [ch02, def:c2-even]; two-component $v^m=x^2b^m-2(b\cdot x)x^m$, $b=v_1$, $\operatorname{div}v=-8\,b\cdot x$ [dictionary, def:gen]. $P_{v_2}$: the translation $x\mapsto v_2$. $\Delta=2E$.
\item $Q_1,S_1,D_1,K_1,P_1,M_1,R_1$: the user's one-loop components along $\xi$, $\eta$, $\Delta$, $K_{v_1}$, translations, rotations of $V$, $\mathfrak{so}(W_{\mathbb R}^\perp)$ (Definition \ref{def:o-comp}).
\item $T\in(0,\infty]$: the endpoint of the transfer, $P(\varepsilon,T)=\int_\varepsilon^T(Q^{GF}\otimes1)K_t\,dt$ [C07, §6.5, Lemma 10.2.1]; for $T=\infty$ this is the propagator of [ch03, def:c3-P].
\end{itemize}
Two-component formulas use Srednicki's contractions $\chi\psi=\chi^\alpha\psi_\alpha$, $\bar\chi\bar\psi=\bar\chi_{\dot\alpha}\bar\psi^{\dot\alpha}$, $\varepsilon^{12}=+1$, the Euclidean $\sigma^m=(\vec\sigma,i)$, $\bar\sigma^m=(-\vec\sigma,i)$, the fields $A_m$, $A_{ij}=A_I\Sigma^I_{ij}$, $A^{ij}=\tfrac12\varepsilon^{ijkl}A_{kl}$, $\lambda_\alpha{}^i$, $\bar\lambda_{\dot\alpha i}$ and the antifields $A^*_m$, $A^*_{ij}$, $\lambda^{*\alpha}{}_i$, $\bar\lambda^{*\dot\alpha i}$, $c^*$ of the dictionary; $(-,-)$ on $\mathfrak g$ contracts the $\mathfrak g$ factors in the order written. Weinberg's form follows from the Weinberg column of the dictionary: $\rho(v)=\gamma_W^\mu v_\mu$ on the Dirac column $(\psi_\alpha,\bar\chi^{\dot\alpha})$ after the Wick map, $\gamma_{5W}=-\gamma_{5S}$, $t_\alpha=gT^a$.
\end{definition}

\begin{lemma}[Rank one]
\label{lem:o-rankone}
For $\xi\in S_+\otimes W$ the ten-component $\Gamma(\xi,\xi)$ vanishes exactly when $\xi$ has rank one, and $[\xi,\xi]=\Gamma(\xi,\xi)|_V=0$ for every $\xi\in S_+\otimes W$ [ch02, prop:c2-special]. For $\xi=u\otimes w$ the image of $\Gamma(\xi,-)|_V$ is two-dimensional.
\end{lemma}

\begin{proof}
By (G1), (G2) of the dictionary,
\[
\Gamma^m(\xi,\xi)\big|_V=2i\,\bar\xi_{\dot\alpha i}\bar\sigma^{m\dot\alpha\alpha}\xi_\alpha{}^i=0,\qquad \Gamma_{kl}(\xi,\xi)=-2\varepsilon_{jikl}\,\xi^{\alpha j}\xi_\alpha{}^i ,
\]
and $\xi^{\alpha j}\xi_\alpha{}^i$ is $\varepsilon^{\alpha\beta}$ contracted with the $2\times2$ minors of the matrix $\xi_\alpha{}^i$; it vanishes for all $i,j$ exactly when the rank is one. For $\xi=u\otimes w$, $\Gamma^m(\xi,\lambda)|_V=i\,\bar\lambda_{\dot\alpha i}w^i\bar\sigma^{m\dot\alpha\alpha}u_\alpha$ ranges over the two-dimensional space $\{u^\alpha\sigma^m_{\alpha\dot\alpha}\bar s^{\dot\alpha}\}$. In the exact realisation of Chapter 2 the weight vector of weights $(1,-1,1,1,1)$ for the Cartan elements $\rho(e_0e_1),\rho(e_2e_3),\rho(e_4e_5),\rho(e_6e_7),\rho(e_8e_9)$ lies in $S_+\otimes W$, has $\Gamma(\xi,\xi)=0$ in all ten components and $\dim\operatorname{Im}\Gamma(\xi,-)|_V=2$.
\end{proof}

\section{The transfer endpoint and the one-loop components}

\begin{definition}[Objects at scale $T$]
\label{def:o-scale}
For $0<\varepsilon<T\le\infty$,
\begin{equation}
P(\varepsilon,T)=\int_\varepsilon^T(Q^{GF}\otimes1)K_t\,dt,\qquad (Q\otimes1+1\otimes Q)P(\varepsilon,T)=K_\varepsilon-K_T\quad[\text{C07, Lemma 10.2.1}],
\label{eq:o-PT}
\end{equation}
with $K_\infty=0$ [ch03, def:c3-K]. Its components are those of [ch03, prop:c3-P] with the scalar factor
\[
\int_\varepsilon^T(4\pi t)^{-2}e^{-|x-y|^2/4t}\,dt=\frac{e^{-|x-y|^2/4T}-e^{-|x-y|^2/4\varepsilon}}{4\pi^2|x-y|^2},\qquad \int_\varepsilon^Te^{-t|p|^2}dt=\frac{e^{-\varepsilon|p|^2}-e^{-T|p|^2}}{|p|^2},
\]
two-component $\langle A_mA_n\rangle$, $\langle A_{ij}A_{kl}\rangle$, $\langle c\,A^*_m\rangle$, $\langle\lambda_\alpha{}^i\bar\lambda_{\dot\beta j}\rangle$ as in Lemma \ref{lem:o-legs} with $e^{-\varepsilon|p|^2}$ replaced by $e^{-\varepsilon|p|^2}-e^{-T|p|^2}$. With the equivariant interaction $I_{\mathfrak{psl}(4|4)}=I+S^{(1)}+S^{(2)}$, its counterterms $I^{CT}_{\mathfrak{psl}(4|4)}[\varepsilon]$, the finite counterterms $F_{\mathfrak{psl}(4|4)}=F_0+F^{(1)}+F^{(2)}$ and $J=I_{\mathfrak{psl}(4|4)}-F_{\mathfrak{psl}(4|4)}-I^{CT}_{\mathfrak{psl}(4|4)}[\varepsilon]+\delta\,O(\varepsilon)$ [ch05, def:c5-Oeps],
\[
I_{\mathfrak{psl}(4|4)}[T]:=\lim_{\varepsilon\to0}W\big(P(\varepsilon,T),J\big)\big|_{\delta=0},\qquad O[T]:=\lim_{\varepsilon\to0}\partial_\delta W\big(P(\varepsilon,T),J\big),
\]
the effective interaction and the observable at scale $T$ [C07], [CG2, §7.2.7], and the scale-$T$ differential on observables [CG2, Def 8.3.1.3]
\begin{equation}
d_T:=d_{CE}+Q+\big\{I_{\mathfrak{psl}(4|4)}[T],-\big\}_T+\hbar\Delta_T,\qquad \Delta_T=-\partial_{K_T}\quad[\text{ch03, def:c3-delta}].
\label{eq:o-dT}
\end{equation}
At $T=\infty$, $\Delta_\infty=0$ and $\{-,-\}_\infty=0$, so $d_\infty=d_{CE}+Q$ and $O[\infty]$ is [ch05, eq:c5-Oinfty].
\end{definition}

\begin{lemma}[The transfer at finite $T$]
\label{lem:o-Tflow}
\begin{itemize}
\item[(i)] $(d_{CE}+Q+\hbar\Delta_T)\,e^{\hbar\partial_{P(\varepsilon,T)}}=e^{\hbar\partial_{P(\varepsilon,T)}}\,(d_{CE}+Q+\hbar\Delta_\varepsilon)$.
\item[(ii)] For $\varepsilon<T<T'\le\infty$: $W(P(\varepsilon,T'),J)=W\big(P(T,T'),W(P(\varepsilon,T),J)\big)$, and $(d_{CE}+Q+\hbar\Delta_{T'})\,e^{\hbar\partial_{P(T,T')}}=e^{\hbar\partial_{P(T,T')}}\,(d_{CE}+Q+\hbar\Delta_T)$.
\item[(iii)] UV-boundary formula at scale $T$:
\[
d_T\,O[T]=\lim_{\varepsilon\to0}\partial_\delta\Big[e^{-W(P(\varepsilon,T),J)/\hbar}\,e^{\hbar\partial_{P(\varepsilon,T)}}\Big(\big((d_{CE}+Q)J+\tfrac12\{J,J\}_\varepsilon+\hbar\Delta_\varepsilon J\big)e^{J/\hbar}\Big)\Big];
\]
the defect vertex $(d_{CE}+Q)J+\tfrac12\{J,J\}_\varepsilon+\hbar\Delta_\varepsilon J$ is the same for every $T$.
\item[(iv)] The singular part $\mathrm{Sing}_\varepsilon$ of every one-loop graph with edges $P(\varepsilon,T)$ is independent of $T\in(0,\infty]$. Hence $I^{CT}_{\mathfrak{psl}(4|4)}[\varepsilon]$, $\partial_\delta I^{CT}_{\mathfrak{psl}(4|4)}[\varepsilon]$ and $O(\varepsilon)$ are the same at every $T$, and so are the finite counterterms fixed by the master equation at the UV boundary.
\end{itemize}
\end{lemma}

\begin{proof}
(i) As in [ch03, thm:c3-uv], with \eqref{eq:o-PT} in place of [ch03, eq:c3-QP]: $\partial_{P(\varepsilon,T)}Q-Q\,\partial_{P(\varepsilon,T)}=\partial_{K_\varepsilon}-\partial_{K_T}=-\Delta_\varepsilon+\Delta_T$; the operators $\partial_\phi$ commute, so $Q\,\partial^n_{P(\varepsilon,T)}=\partial^n_{P(\varepsilon,T)}Q+n\,\partial^{n-1}_{P(\varepsilon,T)}(\Delta_\varepsilon-\Delta_T)$, and summing, $Q\,e^{\hbar\partial_{P(\varepsilon,T)}}=e^{\hbar\partial_{P(\varepsilon,T)}}(Q+\hbar\Delta_\varepsilon)-\hbar\Delta_Te^{\hbar\partial_{P(\varepsilon,T)}}$; $d_{CE}$ commutes with $\partial_{P(\varepsilon,T)}$. (ii) $P(\varepsilon,T')=P(\varepsilon,T)+P(T,T')$ and $e^{\hbar\partial_{\phi_1+\phi_2}}=e^{\hbar\partial_{\phi_2}}e^{\hbar\partial_{\phi_1}}$; the second identity is (i) with $(\varepsilon,T)$ replaced by $(T,T')$. (iii) Apply (i) to $e^{J/\hbar}$; by the product rule of [ch03, def:c3-delta] at scale $T$, $(d_{CE}+Q+\hbar\Delta_T)e^{W/\hbar}=\hbar^{-1}e^{W/\hbar}\big((d_{CE}+Q)W+\tfrac12\{W,W\}_T+\hbar\Delta_TW\big)$ for even $W$, as in [ch03, cor:c3-qme]; apply $\partial_\delta$ and let $\varepsilon\to0$; the $\delta$-linear part of the left side is $d_TO[T]$ with \eqref{eq:o-dT}. (iv) In the Schwinger representation of the proof of [ch04, lem:c4-local], the coefficients of $\varepsilon^{-k}$ and $\log\varepsilon$ come from the region where the sum of the Schwinger times of the cycle is at most a fixed $\tau_0<T$ (the proof of [ch04, lem:c4-local] uses $\tau_0=1$; the singular coefficients do not depend on $\tau_0$); there the integration domain is independent of $T$, and the region where one Schwinger time exceeds $\tau_0$ gives a function with a limit as $\varepsilon\to0$. The finite counterterms are the $\varepsilon\to0$ limits of local defect vertices [ch04, prop:c4-defect1], independent of $T$.
\end{proof}

\begin{definition}[Local representative and one-loop components]
\label{def:o-comp}
Let $O=\int\omega\,f$ be a homogeneous local operator of Weyl weight $w_O$ and degree $|O|$, built in a fixed scheme (finite counterterms $F_{\mathfrak{psl}(4|4)}$, finite parts of $O(\varepsilon)$, gauge-fixing operator $Q^{GF}$). A local representative of its BV transformation at scale $T$ is a cochain $L(O)\in C^\bullet(\mathfrak{psl}(4|4),\mathscr O_l(\mathscr E))$ of degree $|O|+1$, with values $\int\omega\,(\cdots)$ polynomial in the position and in the jets of the fields, such that
\begin{equation}
d_T\,O[T]=\big(L(O)\big)[T],
\label{eq:o-rep}
\end{equation}
with $M\mapsto M[T]$ the construction of Definition \ref{def:o-scale} extended $C^\bullet(\mathfrak{psl}(4|4))$-linearly. Modulo $\hbar^2$,
\begin{equation}
L(O)=s_{\mathfrak{psl}(4|4)}O+\hbar L_1(O),\qquad s_{\mathfrak{psl}(4|4)}:=d_{CE}+Q+\{I_{\mathfrak{psl}(4|4)},-\},
\label{eq:o-L1}
\end{equation}
where $s_{\mathfrak{psl}(4|4)}^2=0$ by [ch02, eq:c2-mc]. The one-loop components are
\[
\begin{aligned}
&Q_1(O):=\hbar L_1(O)^{(1)}(\xi),\qquad S_1(O):=\hbar L_1(O)^{(1)}(\eta),\\
&D_1(O):=\hbar L_1(O)^{(1)}(\Delta),\qquad K_1(O):=\hbar L_1(O)^{(1)}(K_{v_1}),\qquad P_1,M_1,R_1\ \text{likewise},\\
&\hbar L_1(O)^{(2)}(x_1,x_2)\ \text{(arity two)},\qquad \hbar L_1(O)^{(0)}\ \text{(the one-loop BRST transformation)}.
\end{aligned}
\]
The classical parts are, for $|O|=0$ (shuffle sign of [ch02, def:c2-dce] for the $(1,0)$ and $(2,0)$ shuffles),
\[
\begin{aligned}
&(s_{\mathfrak{psl}(4|4)}O)^{(0)}=(Q+\{I,-\})O,\qquad (s_{\mathfrak{psl}(4|4)}O)^{(1)}(x_1)=\{S^{(1)}(x_1),O\},\\
&(s_{\mathfrak{psl}(4|4)}O)^{(2)}(x_1,x_2)=\{S^{(2)}(x_1,x_2),O\}.
\end{aligned}
\]
\end{definition}

The two-component form of the classical part along $\xi=u\otimes w$ is the bracket of $O$ with [dictionary, (F1)], [targets, §2.2]:
\begin{equation}
\begin{aligned}
S^{(1)}(\xi)=\int\Big\{&-i\big(\bar\lambda_{\dot\alpha i}w^i\bar\sigma^{m\dot\alpha\alpha}u_\alpha,\,A^*_m\big)+\tfrac12\varepsilon_{jikl}w^j\big((u\lambda^i),A^{*kl}\big)-\tfrac14\big(F_{mn}(\sigma^m\bar\sigma^nu)_\alpha,\,\lambda^{*\alpha}{}_iw^i\big)\\
&-\tfrac14\big([A^{ik},A_{kj}w^j]u_\alpha,\,\lambda^{*\alpha}{}_i\big)-\tfrac i2\big(u^\alpha\sigma^m_{\alpha\dot\alpha}D_m(A_{ij}w^j),\,\bar\lambda^{*\dot\alpha i}\big)\Big\},
\end{aligned}
\label{eq:o-S1xi}
\end{equation}
index-free $S^{(1)}(\xi)=\int\mathrm{dvol}\big(-(\Gamma(\xi,\lambda),A^*)+\tfrac12(\rho(F_A)\xi,\lambda^*)\big)$ [ch02, eq:c2-S1Q].

\begin{proposition}[Uniqueness]
\label{prop:o-unique}
At every $T\in(0,\infty]$: if a local representative exists, $L_1(O)$ is unique.
\end{proposition}

\begin{proof}
Let $M$ be a local cochain of order $\hbar$ with $M[T]=0$ at order $\hbar$; at this order $M[T]$ is the tree-level flow of $M$ with vertices $I$, $S^{(1)}$, $S^{(2)}$ and edges $P(0,T)$. Induction on the arity $k$: if $M^{(j)}=0$ for $j<k$, the arity-$k$ part of $M[T]$ is the tree-level flow of $M^{(k)}$ with vertices $I$; a tree with $r\ge1$ vertices $I$ has at least $r$ legs more than its vertex $M^{(k)}$, so the part with the fewest legs is the corresponding Taylor component of $M^{(k)}$, which therefore vanishes; by induction on the number of legs, as in [ch04, prop:c4-unique], $M^{(k)}=0$.
\end{proof}

\begin{proposition}[Independence of the transfer endpoint]
\label{prop:o-Tindep}
Assume the one-loop equivariant master equation at the UV boundary (after finite counterterms), so that $d_T^2=0$ for every $T$. If \eqref{eq:o-rep} holds at one $T_0\in(0,\infty)$ with a local $L(O)$, it holds at every $T\in(0,\infty)$ with the same $L(O)$. If moreover the graph sums of $O[T]$ and of $L(O)[T]$ converge as $T\to\infty$, it holds at $T=\infty$ with the same $L(O)$.
\end{proposition}

\begin{proof}
By Lemma \ref{lem:o-Tflow} (ii) at $\delta$-linear order, $O[T]e^{I[T]/\hbar}=e^{\hbar\partial_{P(T_0,T)}}\big(O[T_0]e^{I[T_0]/\hbar}\big)$ for $T>T_0$, with $I[T]:=I_{\mathfrak{psl}(4|4)}[T]$, and the same for $L(O)$; for $T<T_0$ use the inverse operator $e^{-\hbar\partial_{P(T,T_0)}}$. By the master equation at scale $T$, $(d_{CE}+Q+\hbar\Delta_T)\big(O[T]e^{I[T]/\hbar}\big)=\big(d_TO[T]\big)e^{I[T]/\hbar}$. The second identity of Lemma \ref{lem:o-Tflow} (ii) gives $\big(d_TO[T]\big)e^{I[T]/\hbar}=e^{\hbar\partial_{P(T_0,T)}}\big((L(O)[T_0])e^{I[T_0]/\hbar}\big)=\big(L(O)[T]\big)e^{I[T]/\hbar}$. Uniqueness (Proposition \ref{prop:o-unique}) at scale $T$ identifies the local representative. The last statement is the limit $T\to\infty$ of the identity, term by term in the graph expansion.
\end{proof}

\begin{theorem}[The cochain identity]
\label{thm:o-cocycle}
If local representatives exist for $O$ and for the cochain-valued operator $s_{\mathfrak{psl}(4|4)}O$, then, modulo $\hbar^2$,
\begin{equation}
s_{\mathfrak{psl}(4|4)}\,L_1(O)+L_1\big(s_{\mathfrak{psl}(4|4)}O\big)=0 ,
\label{eq:o-cocycle}
\end{equation}
with $L_1$ extended $C^\bullet(\mathfrak{psl}(4|4))$-linearly. In components on two generators $x_1,x_2$, with $|L_1(O)|=|O|+1$, the $d_{CE}$ term of the first summand is $(-1)^{|O|+1}L_1(O)^{(1)}([x_1,x_2])$ [ch02, eq:c2-dce]. The identity involves the local cochains only and is the same at every $T$.
\end{theorem}

\begin{proof}
Apply $d_T$ to \eqref{eq:o-rep}: $0=d_T\big(L(O)[T]\big)=\big(L(L(O))\big)[T]$, and $L(L(O))=s_{\mathfrak{psl}(4|4)}^2O+\hbar\big(s_{\mathfrak{psl}(4|4)}L_1(O)+L_1(s_{\mathfrak{psl}(4|4)}O)\big)+O(\hbar^2)$. By Proposition \ref{prop:o-unique} the order-$\hbar$ coefficient vanishes. For $n=2$ in [ch02, eq:c2-dce] the exponent is $|f|$ modulo $2$.
\end{proof}

\section{The UV-boundary formula and the graph classes}

\begin{proposition}[UV-boundary formula for the components]
\label{prop:o-uv}
With $J$ of Definition \ref{def:o-scale}, $d_TO[T]$ is given by Lemma \ref{lem:o-Tflow} (iii); at $T=\infty$ this is [ch06, eq:c6-uvobs]. By [ch02, eq:c2-mc], for $|O|=0$ and modulo $\hbar^2$ the defect vertex is
\begin{equation}
\begin{aligned}
&\underbrace{\tfrac12\big(\{I_{\mathfrak{psl}(4|4)},I_{\mathfrak{psl}(4|4)}\}_\varepsilon-\{I_{\mathfrak{psl}(4|4)},I_{\mathfrak{psl}(4|4)}\}\big)}_{\text{action special edges}}+\delta\Big[\underbrace{s_{\mathfrak{psl}(4|4)}O}_{\text{classical}}+\underbrace{\{I_{\mathfrak{psl}(4|4)},O\}_\varepsilon-\{I_{\mathfrak{psl}(4|4)},O\}}_{\text{operator special edges}}\Big]\\
&+\hbar\Delta_\varepsilon I_{\mathfrak{psl}(4|4)}-(d_{CE}+Q)\big(I^{CT}_{\mathfrak{psl}(4|4)}[\varepsilon]+F_{\mathfrak{psl}(4|4)}\big)-\big\{I_{\mathfrak{psl}(4|4)},I^{CT}_{\mathfrak{psl}(4|4)}[\varepsilon]+F_{\mathfrak{psl}(4|4)}\big\}_\varepsilon\\
&+\delta\Big[\hbar\Delta_\varepsilon O-(d_{CE}+Q)\,\partial_\delta I^{CT}_{\mathfrak{psl}(4|4)}[\varepsilon]-\big\{I_{\mathfrak{psl}(4|4)},\partial_\delta I^{CT}_{\mathfrak{psl}(4|4)}[\varepsilon]\big\}_\varepsilon-\big\{I^{CT}_{\mathfrak{psl}(4|4)}[\varepsilon]+F_{\mathfrak{psl}(4|4)},O\big\}_\varepsilon\Big].
\end{aligned}
\label{eq:o-defect}
\end{equation}
Every special edge carries $K_\varepsilon-K_0$ [ch04, prop:c4-defect1]; for odd $O$ the signs are those of the proof of [ch06, lem:c6-bvexact]. The defect vertex does not involve $T$; the edges of the graphs are $P(\varepsilon,T)$.
\end{proposition}

\begin{proof}
Lemma \ref{lem:o-Tflow} (iii). Expanding $(d_{CE}+Q)J+\tfrac12\{J,J\}_\varepsilon+\hbar\Delta_\varepsilon J$ in $\hbar$ and $\delta$, the order-$\hbar^0$, $\delta$-free part is $(d_{CE}+Q)I_{\mathfrak{psl}(4|4)}+\tfrac12\{I_{\mathfrak{psl}(4|4)},I_{\mathfrak{psl}(4|4)}\}_\varepsilon$, and [ch02, eq:c2-mc] in the form $(d_{CE}+Q)I_{\mathfrak{psl}(4|4)}=-\tfrac12\{I_{\mathfrak{psl}(4|4)},I_{\mathfrak{psl}(4|4)}\}$ [ESW, Rem 1.29] leaves the special-edge difference; the $\delta$-linear order-$\hbar^0$ part is $(d_{CE}+Q)O+\{I_{\mathfrak{psl}(4|4)},O\}_\varepsilon$; the remaining terms are of order $\hbar$, the self-contraction $\hbar\Delta_\varepsilon$ of the counterterms being of order $\hbar^2$.
\end{proof}

The reduction of the defect by the classical master equation uses the equation as an identity of densities modulo total derivatives, inside graphs whose other vertices are joined to the defect by propagators. At finite $T$ every component of $P(\varepsilon,T)$ and its derivatives decay as $e^{-|y|^2/4T}$ at large $|y|$, so the surface terms vanish for every component of this chapter, including the components along $K_{v_1}$. At $T=\infty$ they vanish when every vertex with position-dependent coefficients decays faster than $|y|^{-4}$ (Theorem \ref{thm:o-ir}), which holds for every component except those whose graphs contain a floating $S^{(1)}(K_{v_1})$ and the pairs $(\eta_1,\eta_2)$ with $\Gamma(\eta_1,\eta_2)|_V\neq0$ (Section 6).

\begin{lemma}[Placement of legs]
\label{lem:o-legs}
$P(\varepsilon,T)$ has the components $A\otimes A$, $c\otimes A^*|_V$, $A^*|_V\otimes c$, $\lambda\otimes\lambda$ only [ch03, prop:c3-P]; two-component [dictionary, prop:prop]
\[
\begin{aligned}
&\langle A_mA_n\rangle:\ 2(-,-)^{-1}\delta_{mn},\qquad \langle A_{ij}A_{kl}\rangle:\ 4(-,-)^{-1}\varepsilon_{ijkl},\\
&\langle c\,A^*_m\rangle:\ -(-,-)^{-1}ip_m,\qquad \langle\lambda_\alpha{}^i\bar\lambda_{\dot\beta j}\rangle:\ -(-,-)^{-1}p_m\sigma^m_{\alpha\dot\beta}\delta^i_j ,
\end{aligned}
\]
each times $(e^{-\varepsilon|p|^2}-e^{-T|p|^2})/|p|^2$. Consequently:
\begin{itemize}
\item a leg $\lambda^*$, $c^*$ or $A^*|_{W_{\mathbb R}^\perp}$ lies on a special edge, on a $K_\varepsilon$ self-contraction, or is external;
\item a vertex $S^{(1)}(x_1)$ with both legs on edges $P(\varepsilon,T)$ (a floating insertion) is, for even $x_1$ with vector field $v$, the term $(v.A|_V,A^*|_V)$ joining an $A|_V$ line to a $c$ line; for odd $x_1$ with Killing spinor $\epsilon(x)$, the term $-(\Gamma(\epsilon(x),\lambda)|_V,A^*|_V)$ joining a $\lambda$ line to a $c$ line; two-component $(v^n\partial_nA_m+A_n\partial_mv^n,A^*_m)$ and $-i\big(\bar\lambda_{\dot\alpha i}\bar\sigma^{m\dot\alpha\alpha}\epsilon_\alpha{}^i+\bar\epsilon_{\dot\alpha i}\bar\sigma^{m\dot\alpha\alpha}\lambda_\alpha{}^i,\,A^*_m\big)$;
\item two floating odd insertions are adjacent only through a $\lambda\otimes\lambda$ edge (a $P$-edge pair); a floating even insertion and a floating odd insertion are never adjacent;
\item $S^{(2)}(x_1,x_2)$, with legs $\lambda^*\lambda^*$ and $A\,c^*$ [ch02, eq:c2-S2], is never floating: one of its legs lies on a special edge or is external.
\end{itemize}
\end{lemma}

\begin{proof}
Read off the propagator table and the field content of \eqref{eq:o-S1xi}, [ch02, eq:c2-S1even, eq:c2-S1S, eq:c2-S2] and (F1)--(F3) of the dictionary. A floating odd insertion has the legs $\lambda$, $A^*|_V$, a floating even insertion the legs $A|_V$, $A^*|_V$; the edges available between their legs are $\lambda\otimes\lambda$ (odd--odd) and none between $\{\lambda,A^*|_V\}$ and $\{A|_V,A^*|_V\}$, since $A\otimes A^*$ and $A^*\otimes A^*$ vanish.
\end{proof}

\begin{theorem}[Graph classes of the arity-one component]
\label{thm:o-classes}
At order $\hbar$, the arity-one component along a generator $x_1$ of a homogeneous local operator $O$ with $|O|=0$ is the $\varepsilon\to0$ limit, at fixed $T$, of the connected graphs with edges $P(\varepsilon,T)$, exactly one marked vertex taken from \eqref{eq:o-defect}, exactly one vertex carrying $x_1$ and exactly one vertex carrying $\delta$, in the following classes.
\begin{itemize}
\item (T) Marked vertex $\{S^{(1)}(x_1),O\}$, other vertices $I$: the graphs of the operator $\{S^{(1)}(x_1),O\}$, i.e.\ trees (the classical component), one-loop graphs, and trees with one vertex $-I^{CT}[\varepsilon]-F_0$.
\item (A) Marked vertex: the special edge between $S^{(1)}(x_1)$ (quadratic or, for odd $x_1$, cubic part) and $O$; other vertices $I$; one loop.
\item (B) Marked vertex: the special edge between $I$ and $S^{(1)}(x_1)$; one vertex $O$; one loop.
\item (C) Marked vertex: the special edge between $I$ and $O$; one vertex $S^{(1)}(x_1)$ on the cycle (floating or adjacent), on a tree, or on an external leg; one loop.
\item (D) Marked vertex: the special edge between two vertices $I$; one vertex $O$, one vertex $S^{(1)}(x_1)$; one loop.
\item (E1) Trees whose marked vertex has order $\hbar$: $-Q\,\partial_\delta I^{CT}_{\mathfrak{psl}(4|4)}[\varepsilon](x_1)-\{I,\partial_\delta I^{CT}_{\mathfrak{psl}(4|4)}[\varepsilon](x_1)\}_\varepsilon-\{S^{(1)}(x_1),\partial_\delta I^{CT}[\varepsilon]\}_\varepsilon$; $-\{I^{CT}_{\mathfrak{psl}(4|4)}[\varepsilon](x_1)+F^{(1)}(x_1),O\}_\varepsilon$; $-\{I^{CT}[\varepsilon]+F_0,O\}_\varepsilon$ with a vertex $S^{(1)}(x_1)$; the arity-one action defect $-QI^{CT}_{\mathfrak{psl}(4|4)}[\varepsilon](x_1)-\{I,I^{CT}_{\mathfrak{psl}(4|4)}[\varepsilon](x_1)\}_\varepsilon-\{S^{(1)}(x_1),I^{CT}[\varepsilon]\}_\varepsilon-(Q+\{I,-\})F^{(1)}(x_1)-\{S^{(1)}(x_1),F_0\}$ with a vertex $O$; the arity-zero action defect of order $\hbar$ with $O$ and $S^{(1)}(x_1)$; the arity-zero operator defect of order $\hbar$ with $S^{(1)}(x_1)$.
\item (E2) Trees whose marked vertex is one of the order-$\hbar^0$ special edges of (A)--(D), with one vertex $-I^{CT}_{\mathfrak{psl}(4|4)}[\varepsilon]$, $-F_{\mathfrak{psl}(4|4)}$ or $-\partial_\delta I^{CT}_{\mathfrak{psl}(4|4)}[\varepsilon]$ of the complementary arity.
\end{itemize}
A special edge on a bridge carries a generic leg momentum $p$ and the factor $e^{-\varepsilon|p|^2}-1$; it is paired with the $K_\varepsilon-K_0$ part of the corresponding counterterm bracket of class (E2), and the pair tends to $(e^{-\varepsilon|p|^2}-1)$ times a renormalised subgraph, hence to $0$ (Part C, classes I-b and III). Writing $\delta_{x_1}O:=\{S^{(1)}(x_1),O\}$,
\begin{equation}
\hbar L_1(O)^{(1)}(x_1)\big[T\big]_{\rm tree}=\lim_{\varepsilon\to0}\Big[\sum_{(A)\text{--}(E2)}\text{graphs}+\big(\partial_\delta I^{CT}[\varepsilon]\ \text{of the operator }\delta_{x_1}O\big)\big[T\big]_{\rm tree}\Big],
\label{eq:o-L1graphs}
\end{equation}
and $\hbar L_1(O)^{(1)}(x_1)$ is read off from the parts with the fewest legs (Proposition \ref{prop:o-unique}).
\end{theorem}

\begin{proof}
Order $\hbar$ forces $b_1=1$ with all vertices of order $\hbar^0$, or $b_1=0$ with one vertex of order $\hbar$ [ch03, prop:c3-graphs]. The arity-one, $\delta$-linear part of the product of cochains places $x_1$ and $\delta$ on exactly one vertex each. The order-$\hbar^0$ marked vertices of \eqref{eq:o-defect} at arity one, $\delta$-linear, are: $\delta_{x_1}O$; the special edges $S^{(1)}(x_1)$--$O$, $I$--$O$ (with $S^{(1)}(x_1)$ elsewhere), $I$--$S^{(1)}(x_1)$ (with $O$ elsewhere), $I$--$I$ (with both elsewhere): classes (T)--(D). The order-$\hbar$ marked vertices give (E1), the order-$\hbar^0$ special edges with one counterterm vertex give (E2). The class (T) graphs are the graphs of $(\delta_{x_1}O)[T]$ before its counterterm; adding and subtracting $\partial_\delta I^{CT}[\varepsilon](\delta_{x_1}O)$ gives $(\delta_{x_1}O)[T]$ plus \eqref{eq:o-L1graphs}, and $(s_{\mathfrak{psl}(4|4)}O+\hbar L_1(O))[T]$ at arity one equals $(\delta_{x_1}O)[T]+\hbar L_1(O)^{(1)}(x_1)[T]_{\rm tree}$, the trees with $(s_{\mathfrak{psl}(4|4)}O)^{(0)}=(Q+\{I,-\})O$ being of order $\hbar^0$ in class (T) when $O$ is gauge variant.
\end{proof}

\begin{remark}[Completion of the list in the road ahead]
\label{rem:o-roadahead}
[ch07, The road ahead] lists the graphs with one vertex given by the arity-one part of [ch07, eq:c7-defect]; these are classes (A), (B) and the corresponding parts of (E1), (E2). The flow of the defect also contains the arity-zero special edges with a separate vertex $S^{(1)}(x_1)$: classes (C) and (D). For the action alone (no $O$) the same holds: the arity-one obstruction along $x_1$ contains the cycles with an $I$--$I$ special edge and $S^{(1)}(x_1)$ on the cycle.
\end{remark}

\begin{proposition}[Arity two along $(\xi,\xi)$]
\label{prop:o-arity2}
For $\xi=u\otimes w$, with $[\xi,\xi]=0$, the arity-two marked vertices of order $\hbar^0$ are $\{S^{(1)}(\xi),S^{(1)}(\xi)\}_\varepsilon-\{S^{(1)}(\xi),S^{(1)}(\xi)\}$, $\{I,S^{(2)}(\xi,\xi)\}_\varepsilon-\{I,S^{(2)}(\xi,\xi)\}$, $\delta\{S^{(2)}(\xi,\xi),O\}$ and $\delta(\{S^{(2)}(\xi,\xi),O\}_\varepsilon-\{S^{(2)}(\xi,\xi),O\})$; the arity-one marked vertices enter with one further vertex $S^{(1)}(\xi)$, the arity-zero marked vertices with two further vertices $S^{(1)}(\xi)$. Since
\begin{equation}
S^{(2)}(\xi,\xi)=-\tfrac12\int\mathrm{dvol}\,\big((\xi,\lambda^*),(\xi,\lambda^*)\big)=-\tfrac12\int\big(u_\alpha w^i\lambda^{*\alpha}{}_i,\;u_\beta w^j\lambda^{*\beta}{}_j\big),
\label{eq:o-S2xi}
\end{equation}
a vertex $S^{(2)}(\xi,\xi)$ occurs only inside the marked vertices listed, the graphs containing it have an external $\lambda^*$, and the outputs without antifields come from graphs with two vertices $S^{(1)}(\xi)$.
\end{proposition}

\begin{proof}
\eqref{eq:o-S2xi} is [ch02, eq:c2-S2] with $\Gamma(\xi,\xi)=0$ (Lemma \ref{lem:o-rankone}) and (P1) of the dictionary, $(\xi,\lambda^*)=\xi_\alpha{}^i\lambda^{*\alpha}{}_i-\bar\xi_{\dot\alpha i}\bar\lambda^{*\dot\alpha i}$. A graph has one field--antifield pairing (the marked special edge or bracket); both legs of $S^{(2)}(\xi,\xi)$ are $\lambda^*$, which lie on no edge $P(\varepsilon,T)$ (Lemma \ref{lem:o-legs}); a separate vertex $S^{(2)}(\xi,\xi)$ next to a marked vertex of arity zero would therefore be disconnected from it, and inside a marked vertex at most one of its legs is on the pairing, so the other is external [targets, Lemma lem:S2].
\end{proof}

\begin{tabular}{l|l|c|l|l|l}
target & operator (index-free; two-component) & $(w_O,|O|)$ & generator, arity & output legs & allowed output structures \\
\hline
Konishi descendant & $\int\omega\,(\lambda,\rho(A|_{W_{\mathbb R}^\perp})(u'\otimes w))$; $\int\omega\,(u'^\alpha\lambda_\alpha{}^i,A_{ij}w^j)$ & $(\tfrac52,0)$ & $\xi$, 1 & $(\bar\lambda,\bar\lambda)$; $(A,A,A)$ & $B_w$, $\bar F_w$ (Prop \ref{prop:o-konishi}) \\
traceless bilinear, plane of $w$ & $\int\omega\,(\iota_{w_1}A,\iota_{w_2}A)$, $w_{a\,ij}w^j=0$; $\tfrac1{16}\int\omega\,(w_1^{ij}A_{ij},w_2^{kl}A_{kl})$ & $(2,0)$ & $\xi$, 1 & $(A,\lambda)$ & none (Prop \ref{prop:o-planevanish}) \\
Konishi primary & $\int\omega\,(A|_{W_{\mathbb R}^\perp},A|_{W_{\mathbb R}^\perp})$; $\tfrac14\int\omega\,(A^{ij},A_{ij})$ & $(2,0)$ & $\xi\in\Sigma$, 1 & $(A,\lambda)$, $(c,\lambda^*)$ & three [ch07, cor:c7-Q] \\
letters & $\int\omega\,(X,\iota_{w_1}A)$, $\int\omega\,(X,\iota_vA|_V)$, $\int\omega\,(\zeta,\lambda)$, $\int\omega\,(X,c)$ & $(1,0),(1,0),(\tfrac32,0),(0,1)$ & $\xi$, 0 and 1 & one and two legs & weight $w_O+\tfrac12$ \\
$bc$ entry & $\int\omega\,(X,u'^\alpha F_{mn}(\sigma^m\bar\sigma^n)_\alpha{}^\beta u'_\beta)(Y,c)$ & $(2,1)$ & $\xi$, 0, 1, 2 & $(\bar\lambda,\bar\lambda)$ & [targets, T4] \\
Konishi descendants of weight 3 & $B_w$, $\bar F_w$ & $(3,0)$ & $\eta$, 1 & $(\lambda,A)$ & term in $4\eta$, weight $\tfrac52$ \\
scalar bilinears & as above & $(2,0)$ & $K_{v_1}$, 1 & $(A,A)$ & $\tfrac14\operatorname{div}v$ term [ch07, cor:c7-K] \\
\end{tabular}
\tabcaption{Targets of the first engine computations, with the output structures allowed by [ch07, prop:c7-select]. $B_w$, $\bar F_w$ are defined in \eqref{eq:o-BF}.}

\section{The twisted $Q_1$ as a consequence}

\begin{lemma}[Evaluation along $\xi$]
\label{lem:o-ev}
Let $\xi\in\Sigma$ satisfy $[\xi,\xi]=0$ and let $q$ be a variable of degree $1$ and even total parity. The map
\begin{equation}
r_\xi\colon C^\bullet(\mathfrak{psl}(4|4),\mathscr O_l(\mathscr E))\to\mathscr O_l(\mathscr E)[q],\qquad f\mapsto\sum_{k\ge0}\frac{(-1)^{k(k-1)/2}}{k!}\,f^{(k)}(\xi,\dots,\xi)\,q^k,
\label{eq:o-ev}
\end{equation}
intertwines $d_{CE}$ with $0$ and the cochain bracket of [ch02, def:c2-dce] with the $q$-linear extension of $\{-,-\}$. For $\xi=u\otimes w$,
\begin{equation}
S^\xi:=r_\xi\big(S_{\mathfrak{psl}(4|4)}\big)\big|_{q=1}=S+S^{(1)}(\xi)-\tfrac12S^{(2)}(\xi,\xi)=S+S^{(1)}(\xi)+\tfrac14\int\big(u_\alpha w^i\lambda^{*\alpha}{}_i,\;u_\beta w^j\lambda^{*\beta}{}_j\big)
\label{eq:o-twistaction}
\end{equation}
satisfies $\{S^\xi,S^\xi\}=0$, and $S+S^{(1)}(\xi)+c\,S^{(2)}(\xi,\xi)$ satisfies it exactly for $c=-\tfrac12$. The same holds for every $x\in\Sigma$ with $[x,x]=0$ and $\Gamma(x,x)=0$, for instance $x=\xi+\bar\xi'$ with $\bar\xi'\in S_-\otimes W^*$ of rank one and $\bar\xi'_{\dot\alpha i}w^i=0$, except for the statement that $\{S^{(1)}(x),S^{(1)}(x)\}\neq0$, which is checked separately.
\end{lemma}

\begin{proof}
For a $(k,l)$-shuffle $\sigma$ with all arguments equal to the odd $\xi$, $\operatorname{sgn}(\sigma)$ times the Koszul sign is $+1$, and the exponent of [ch02, def:c2-dce] is $|g|k+k(l+|g|)\equiv kl$; hence $\{f,g\}(\xi^{k+l})=\binom{k+l}{k}(-1)^{kl}\{f(\xi^k),g(\xi^l)\}$, and with $a_k=(-1)^{k(k-1)/2}/k!$,
\[
a_ka_l=a_{k+l}\binom{k+l}{k}(-1)^{kl},\qquad\text{since}\quad \tfrac12(k+l)(k+l-1)=\tfrac12k(k-1)+\tfrac12l(l-1)+kl .
\]
$d_{CE}$ evaluated on $(\xi,\dots,\xi)$ contains $[\xi,\xi]=0$. For the twisted action: by [ch02, lem:c2-components] in arity one, two and three with $[\xi,\xi]=0$ and $S^{(3)}=0$, and [ch02, prop:c2-tower],
\[
\begin{aligned}
&\{S,S^{(1)}(\xi)\}=0,\qquad \{S,S^{(2)}(\xi,\xi)\}=\{S^{(1)}(\xi),S^{(1)}(\xi)\},\\
&\{S^{(1)}(\xi),S^{(2)}(\xi,\xi)\}=0,\qquad \{S^{(2)},S^{(2)}\}=0,
\end{aligned}
\]
so $\{S+S^{(1)}(\xi)+cS^{(2)}(\xi,\xi),\ \cdot\ \}$ squared gives $(2c+1)\{S^{(1)}(\xi),S^{(1)}(\xi)\}$. This functional is nonzero: it equals $\{S,S^{(2)}(\xi,\xi)\}$, whose $\lambda^*\lambda$ part is $-\int\mathrm{dvol}\,((\xi,\lambda^*),(\xi,{\not d}_A\lambda))$ up to the sign of the pair, and $(\xi,{\not d}\lambda)=i\,u_\alpha w^i\,\partial_m\bar\lambda_{\dot\alpha i}\bar\sigma^{m\dot\alpha\alpha}\neq0$ by (R1). In the exact realisation of Chapter 2 with $\mathfrak g=\mathfrak{su}(2)$ and the rank-one $\xi$ of Lemma \ref{lem:o-rankone}, $\{S^{(1)}(\xi),S^{(1)}(\xi)\}$ is a nonzero local functional, $\{S^{(1)}(\xi),S^{(1)}(\xi)\}-\{S,S^{(2)}(\xi,\xi)\}=0$ and $\{S^{(1)}(\xi),S^{(2)}(\xi,\xi)\}=0$. The two-component form uses \eqref{eq:o-S2xi}. For $x=\xi+\bar\xi'$: $\Gamma(x,x)|_V=2\Gamma(\xi,\bar\xi')|_V$ is proportional to $\bar\xi'_{\dot\alpha i}w^i=0$ by (G1), and the components of $\Gamma$ along $W_{\mathbb R}^\perp$ pair equal chiralities and vanish on the rank-one $\xi$, $\bar\xi'$ by (G2); the computation above uses only $[x,x]=0$ and $S^{(3)}=0$.
\end{proof}

\begin{lemma}[No action-side cochains of arity three or more along $\Sigma$]
\label{lem:o-actiontrunc}
Write $m(n)$ for the smallest Weyl weight of a local functional of degree $-n$: $m(n)=2n$ for $n$ even and $m(n)=2n+\tfrac12$ for $n$ odd ($n\ge1$). The arity-$k$ value on $\Sigma^{\otimes k}$ of an equivariant counterterm or finite counterterm of degree $0$ has degree $-k$ and weight $4+\tfrac k2$; since $m(k)>4+\tfrac k2$ for $k\ge3$, it vanishes for $k\ge3$. In particular $r_\xi$ of the action-side cochains $I^{CT}_{\mathfrak{psl}(4|4)}[\varepsilon]$, $F_{\mathfrak{psl}(4|4)}$ stops at $q^2$.
\end{lemma}

\begin{proof}
Degrees and weights of the fields: $c$: $(+1,0)$, $A$: $(0,1)$, $\lambda$: $(0,\tfrac32)$, $A^*$: $(-1,3)$, $\lambda^*$: $(-1,\tfrac52)$, $c^*$: $(-2,4)$, a derivative $(0,1)$ [ch07, §Allowed output structures]. A factor $c$ raises the degree at weight $0$, so the smallest weight at degree $-n$ uses $c^*$ and at most one $\lambda^*$: $2n$, resp.\ $2(n-1)+\tfrac52$. The weight of the value: the cochain has weight $4$, each argument in $\Sigma$ carries weight $-\tfrac12$ [ch07, prop:c7-select]. For $k=3$: $m(3)=\tfrac{13}{2}>\tfrac{11}2$; for $k\ge4$, $m(k)\ge2k>4+\tfrac k2$. The coefficients $\varepsilon^{-1}$ of the counterterms carry weight $2$ and only lower the weight left to the fields. Machine enumeration of all field contents with at most six factors confirms $m(n)$ for $n\le4$ and the absence of contents at $k=3,4,5$.
\end{proof}

\begin{remark}[The coefficient of the arity-two term in the twisted action]
[ESW, Def 2.17] prints $S^Q=S+\sum_{k\ge1}S^{(k)}_{\mathcal A}(Q,\dots,Q)$. With [ESW, (34)] and [ESW, Lemma 3.6] as reproduced in [ch02, lem:c2-dcesign], the coefficient of $S^{(2)}(\xi,\xi)$ for which the twisted action solves the classical master equation is $-\tfrac12$ (Lemma \ref{lem:o-ev}); $-\tfrac12=a_2$ is the décalage coefficient of \eqref{eq:o-ev}. This settles open item 2 of the targets file for the action: in $\mathrm{ev}_\xi=\sum_kc_kF^{(k)}(\xi,\dots,\xi)$, $c_0=c_1=1$, $c_2=-\tfrac12$. For the one-loop components the same coefficients apply under hypothesis (H1) of Theorem \ref{thm:o-twist}.
\end{remark}

\begin{definition}[Twisted differentials]
\label{def:o-twisted}
For $\xi=u\otimes w$ (or any $x$ as in Lemma \ref{lem:o-ev}):
\[
\begin{aligned}
&Q_0^\xi:=\{S^\xi,-\}=Q+\big\{I+S^{(1)}(\xi)-\tfrac12S^{(2)}(\xi,\xi),\,-\big\},\\
&Q_1^\xi(O):=r_\xi\big(\hbar L_1(O)\big)\big|_{q=1}=\hbar\sum_{k\ge0}\frac{(-1)^{k(k-1)/2}}{k!}\,L_1(O)^{(k)}(\xi,\dots,\xi).
\end{aligned}
\]
In the names of Definition \ref{def:o-comp} the terms $k=0,1$ are $\hbar L_1(O)^{(0)}+Q_1(O)$.
\end{definition}

\begin{proposition}[Truncation for the targets]
\label{prop:o-trunc}
The arity-$k$ value $L_1(O)^{(k)}(\xi,\dots,\xi)$ has Weyl weight $w_O+\tfrac k2$ and degree $|O|+1-k$; it vanishes unless $w_O+\tfrac k2\ge m(k-1-|O|)$ (Lemma \ref{lem:o-actiontrunc}), which implies $k\le\tfrac13(2w_O+4+4|O|)$. For the operators of the target table and $\xi=u\otimes w$, $L_1(O)^{(k)}(\xi,\dots,\xi)=0$ for $k\ge3$, so
\[
Q_1^\xi(O)=\hbar\Big(L_1(O)^{(0)}+L_1(O)^{(1)}(\xi)-\tfrac12L_1(O)^{(2)}(\xi,\xi)\Big).
\]
For operators of higher weight the census of the arity-$k$ components with $k\ge3$ allowed by the bound is part of the computation.
\end{proposition}

\begin{proof}
[targets, Lemma lem:arity]: each argument carries weight $-\tfrac12$ and lowers the degree by one. Letters ($w_O\le\tfrac32$, $|O|=0$; $w_O=0$, $|O|=1$) and scalar bilinears ($w_O=2$, $|O|=0$): the bound excludes $k\ge3$. Konishi descendant $O_\zeta$ ($w_O=\tfrac52$): $k=3$ requires weight $4$ and degree $-2$, the content $c^*$ alone, which is not a colour singlet; besides, the value carries $w^{\otimes4}$ ($w^{\otimes3}$ from the arguments, one $w$ from $\zeta=u'\otimes w$) and $c^*$ has no $W$ index; $k\ge4$ is excluded by the bound. $bc$ entry ($w_O=2$, $|O|=1$): $k=3$ requires weight $\tfrac72$ and degree $-1$, the contents $(\lambda^*,A)$ and $\partial\lambda^*$; the value carries $w^{\otimes3}$, which needs three lower $W$ indices not two of which sit in one $A_{ij}$, since $A_{ij}w^iw^j=0$; $\lambda^*$ has one lower index and $A_{ij}$ an antisymmetric pair, so both contents give $0$; $k=4$ requires weight $4$ and degree $-2$, the content $c^*$, with no $W$ index; $k\ge5$ is excluded by the bound. Weight-$3$ operators $B_w$, $\bar F_w$ along $\xi$: $k=3$ requires weight $\tfrac92$ at degree $-2$, and no content exists. The enumeration is machine-checked.
\end{proof}

\begin{theorem}[Recovery of the twisted $Q_1$]
\label{thm:o-twist}
Let $\xi=u\otimes w$ (or $x$ as in Lemma \ref{lem:o-ev}). Assume:
\begin{itemize}
\item[(H1)] the products of cochain-valued functionals in $W(P,J)$ and the $C^\bullet(\mathfrak{psl}(4|4))$-linear extension of $L_1$ carry the décalage signs for which $r_\xi$ is multiplicative; these signs are fixed in code (engine E0) before the coefficient $-\tfrac12$ of $L_1^{(2)}(\xi,\xi)$ is quoted;
\item[(H2)] the one-loop equivariant master equation holds in arities $0,1,2$ along $\xi$ after finite counterterms (arity $\ge3$ along $\Sigma$ is empty by Lemma \ref{lem:o-actiontrunc});
\item[(H3)] local representatives exist for the operators considered; for (ii) and (iii) on all of $\mathscr O_l$.
\end{itemize}
Then:
\begin{itemize}
\item[(i)] $Q\big(r_\xi O[\infty]|_{q=1}\big)=\big(Q_0^\xi O+Q_1^\xi(O)\big)[\infty]_\xi+O(\hbar^2)$, where $[\infty]_\xi$ is the construction [ch05, eq:c5-Oinfty] with the interaction $r_\xi(I_{\mathfrak{psl}(4|4)}-F_{\mathfrak{psl}(4|4)}-I^{CT}_{\mathfrak{psl}(4|4)}[\varepsilon])|_{q=1}$; this is the twisted quantum differential on twisted observables at scale $\infty$, where $\Delta_\infty=0$ [ch03, def:c3-delta]. At finite $T$ the same statement holds with $Q$ replaced by $Q+\{r_\xi I_{\mathfrak{psl}(4|4)}[T]|_{q=1},-\}_T+\hbar\Delta_T$ and $[\infty]_\xi$ by $[T]_\xi$; the limit $T\to\infty$ is safe along $\xi$ (Theorem \ref{thm:o-ir} (i), (iv)).
\item[(ii)] $(Q_0^\xi)^2=0$ and $Q_0^\xi\,Q_1^\xi+Q_1^\xi\,Q_0^\xi=0$ on local functionals.
\item[(iii)] The induced map $[Q_1^\xi]\colon H(\mathscr O_l,Q_0^\xi)\to H(\mathscr O_l,Q_0^\xi)$ is invariant under the scheme changes (a), (b), (d) with $G$ the coupling direction, and (e) of Proposition \ref{prop:o-laws}; under (d) for every $s_{\mathfrak{psl}(4|4)}$-closed $G$ when moreover (H4) holds: every $s_{\mathfrak{psl}(4|4)}$-closed local cochain of degree $0$ and weight $4$ lies in $\mathbb C(S-S^{(2)})\oplus\mathbb C\,G_B+\operatorname{Im}s_{\mathfrak{psl}(4|4)}$ (engine E3b). Item (c) of Proposition \ref{prop:o-laws} is a proof plan and enters no statement.
\item[(iv)] Two one-loop quantizations of the same twisted classical theory whose difference is the $Q_0^\xi$-closed degree-$0$ local functional $D$ have class maps differing by $[\{D,-\}]$. If $D\in\mathbb C\,S^\xi+\mathbb C\,N_{R_\xi}S^\xi+\operatorname{Im}Q_0^\xi$ (with $N_{R_\xi}$ of Proposition \ref{prop:o-laws}), the class maps coincide.
\end{itemize}
\end{theorem}

\begin{proof}
(i) By Lemma \ref{lem:o-ev}, $r_\xi$ intertwines $d_{CE}+Q$ with $Q$ and the bracket with the bracket; by (H1) it is multiplicative on the products of the flow, which act on the field dependence of the values and commute with the evaluation of the cochain arguments. Apply $r_\xi$ to \eqref{eq:o-rep} with \eqref{eq:o-L1}; at finite $T$ apply it to \eqref{eq:o-rep} with \eqref{eq:o-dT}. (ii) $(Q_0^\xi)^2=\tfrac12\{\{S^\xi,S^\xi\},-\}=0$ by Lemma \ref{lem:o-ev}; the second identity is $r_\xi$ of \eqref{eq:o-cocycle}, which by (H1), (H3) and Proposition \ref{prop:o-trunc} for the higher components is an identity of finite sums. (iii) By \eqref{eq:o-twistlaws} each listed change adds to $Q_1^\xi$ a term $aQ_0^\xi$ or a graded commutator $[Q_0^\xi,\Psi]$, and both induce $0$ on $H(Q_0^\xi)$: $\pi_0Q_0^\xi=0$, $Q_0^\xi\iota_0=0$ [targets, Lemma lem:inv]. Under (H4) a general closed $G$ is a combination of the coupling direction, $G_B$ and $s\chi$, and $r_\xi(s\chi)=Q_0^\xi r_\xi(\chi)$ gives $-\{Q_0^\xi r_\xi\chi,-\}=-[Q_0^\xi,\{r_\xi\chi,-\}]$. (iv) Changing the order-$\hbar$ interaction by $-\hbar D$ changes $\hbar L_1$ by $-\hbar\{D,-\}$ (proof of Proposition \ref{prop:o-laws}, item (d)); $\{D,-\}$ commutes with $Q_0^\xi$ up to sign because $Q_0^\xi D=0$; for $D=S^\xi$ it is $Q_0^\xi$, for $D=N_{R_\xi}S^\xi$ it is $[N_{R_\xi},Q_0^\xi]$, both zero on $H(Q_0^\xi)$.
\end{proof}

\begin{remark}[Transport to the holomorphic letters: an identification]
\label{rem:o-transport}
[ESW, Thm 10.9]: the holomorphic twist on $\mathbb R^4$ is perturbatively equivalent to the BF theory with fields $T^*[-1]\mathrm{Map}(\mathbb C^2_{Dol},B\mathfrak g)$, $U(2)$-equivariantly. Its proof organises the fields, in ESW's own symbols, into the $U(2)$-equivariant collections listed after it, and eliminates trivial pairs [ESW, proof of Thm 10.1, Props 1.21, 1.23]. The quasi-isomorphism $\Phi$ of local-operator complexes that transports $[Q_1^\xi]$ to the $Q_0$-cohomology of the holomorphic letters $c,\gamma,\beta,b$ of [BK] is to be derived from these collections, with the regrading to [BK]'s letters written out, before any numerical comparison. The letter table of [targets, §2.3] is an identification at the level of free $Q_0^\xi$-cohomology; it is used for comparison only.
\end{remark}

\begin{remark}[Comparison with the twisted computation]
[BK, (2.42), footnote 22] and [W18, Thm 3.4] quantize the holomorphic theory with the holomorphic gauge fixing and no counterterms; the restriction $r_\xi$ of this chapter uses the untwisted Feynman-gauge propagator and treats $S^{(1)}(\xi)$ as a vertex; the commutator of the linear part of $S^{(1)}(\xi)$ with $Q^{GF}$ is nonzero [ch07, prop:c7-susy], and its effect is carried by the equivariant master equation along $\xi$. By Theorem \ref{thm:o-twist} (iv), agreement at the level of classes holds when the difference of the two quantizations lies in $\mathbb C\,S^\xi+\mathbb C\,N_{R_\xi}S^\xi+\operatorname{Im}Q_0^\xi$; [BK] state that the twisted action carries no coupling [BK, (1.3), (2.22)]. The engine checks this on the classes of [targets, T3--T5, T12] (Section 7, E1, E4, E7, E12).
\end{remark}

\begin{remark}[Reduction of outputs]
\label{rem:o-exact}
Only full $s_{\mathfrak{psl}(4|4)}$-exact cochains are discarded: for a local cochain $\chi$, $r_\xi(s_{\mathfrak{psl}(4|4)}\chi)=Q_0^\xi\,r_\xi(\chi)$, with $\chi^{(0)}$, $\chi^{(1)}(\xi)$, $\chi^{(2)}(\xi,\xi)$ all carried along. A term $(Q+\{I,-\})\chi^{(1)}(\xi)$ dropped at arity one alone differs from $Q_0^\xi\chi^{(1)}(\xi)$ by $\{S^{(1)}(\xi),\chi^{(1)}(\xi)\}$ and by $S^{(2)}$ terms of higher arity, which then have to be dropped at arity two with the same $\chi$. The engines record each reduction as the triple $(\chi^{(0)},\chi^{(1)},\chi^{(2)})$.
\end{remark}

\begin{tabular}{l|l|l}
target & statement in this chapter & status \\
\hline
T1 & Theorem \ref{thm:o-twist} (ii); $[Q_1^\xi]^2=0$ on $H(Q_0^\xi)$ needs order $\hbar^2$ & proved under (H1)--(H3) / Step 5 \\
T2 & pair rule and single letters: families of Proposition \ref{prop:o-laws} & engine E6 \\
T3 & Proposition \ref{prop:o-planevanish} & proved under [ch07, prop:c7-select] \\
T4 & arity two along $(\xi,\xi)$, Proposition \ref{prop:o-arity2} & engine E7 \\
T5 & Proposition \ref{prop:o-konishi}: $c_K$ is a scheme-dependent pair value & engine E1, E2 \\
T6 & arity one with field representatives & engine E7 \\
T7, T8, T9, T11 & classes, Theorem \ref{thm:o-invariants} & from E1--E7 \\
T10 & $(\xi,\eta)$ component with $[\xi,\eta]=0$, Proposition \ref{prop:o-S} (iii) & engine E8 \\
T12 & rank-$(1,1)$ element $x=\xi+\bar\xi'$: $[Q_1^x]=0$ on $H(Q_0^x)$ & engine E12 \\
\end{tabular}
\tabcaption{The targets of [targets, Section 3] and where they are reproduced.}

\section{Renormalisation freedom}

\begin{definition}[The outer derivation]
\label{def:o-B}
$B$ is the even linear map of $\mathfrak{psl}(4|4)$ equal to $0$ on the even part, to $+1$ on $S_+\otimes W\subset\Sigma$ and on $S_-\otimes W\subset\Sigma^*$, and to $-1$ on $S_-\otimes W^*\subset\Sigma$ and on $S_+\otimes W^*\subset\Sigma^*$. Two-component: $B$ multiplies $\xi_\alpha{}^i$ and $\bar\eta^{\dot\alpha i}$ by $+1$ and $\bar\xi_{\dot\alpha i}$ and $\eta^\alpha{}_i$ by $-1$, i.e.\ it counts upper minus lower $W$ indices of the odd parameters; it is induced by $\mathrm{ad}$ of $\mathrm{diag}(1_4,0_4)$ in $\mathfrak{gl}(4|4)$.
\end{definition}

\begin{lemma}[$B$ is an outer derivation]
\label{lem:o-Bouter}
$B$ is an even derivation of $\mathfrak{psl}(4|4)$, and $B\neq\mathrm{ad}_y$ for every $y\in\mathfrak{psl}(4|4)$.
\end{lemma}

\begin{proof}
On odd $x,y$ with $B$-values $b_x,b_y$ the derivation rule reads $(b_x+b_y)[x,y]=0$, since $[x,y]$ is even. Brackets within one $B$-eigenspace vanish: $[S_+\otimes W,S_+\otimes W]$ and $[S_-\otimes W^*,S_-\otimes W^*]$ lie in the translations $V\cong S_+\otimes S_-$ and have the Lorentz type $S_\pm\otimes S_\pm$; $[S_+\otimes W,S_-\otimes W]$ and $[S_-\otimes W^*,S_+\otimes W^*]$ have the Lorentz type $S_+\otimes S_-$, which does not occur in $[\Sigma,\Sigma^*]\subset\mathfrak{so}(V)\oplus\mathfrak{so}(W_{\mathbb R}^\perp)\oplus\mathbb C\Delta$; $[S_+\otimes W^*,S_+\otimes W^*]$ and $[S_-\otimes W,S_-\otimes W]$ vanish by $\Gamma|_V$ [ch02, prop:c2-special]. The even part preserves the eigenspaces: translations map $\Sigma^*\to\Sigma$ by $\rho(v_2)$, special conformal fields $\Sigma\to\Sigma^*$ by $\rho(x)$, both keeping the $W$ factor, and $\mathfrak{so}(V)$, $\Delta$, $\mathfrak{so}(W_{\mathbb R}^\perp)$ preserve each summand. In the exact realisation of Chapter 2 the four chirality blocks, the vanishing of all brackets within the two $B$-eigenspaces and the preservation by $P,K,\Delta,\mathfrak{so}(V),\mathfrak{so}(W_{\mathbb R}^\perp)$ are machine-checked. If $B=\mathrm{ad}_y$, then $y$ centralises the even part $\cong\mathfrak{sl}(4)\oplus\mathfrak{sl}(4)$, which is semisimple, and $\Sigma\oplus\Sigma^*$ contains no invariant vector; hence $y=0$, while $B\neq0$.
\end{proof}

\begin{proposition}[Transformation laws]
\label{prop:o-laws}
Keep the labels $O$ of the operators fixed. Under the following changes of scheme, $L_1$ changes as stated ($\mathrm{ad}_\phi=\{\phi,-\}$; the shuffle signs of [ch02, def:c2-dce]).
\begin{itemize}
\item[(a)] Exact finite counterterm: $F_{\mathfrak{psl}(4|4)}\mapsto F_{\mathfrak{psl}(4|4)}+\hbar\,s_{\mathfrak{psl}(4|4)}\phi$, $\phi$ a local cochain of degree $-1$ (a field redefinition): $L_1\mapsto L_1+[s_{\mathfrak{psl}(4|4)},\mathrm{ad}_\phi]=L_1-\mathrm{ad}_{s_{\mathfrak{psl}(4|4)}\phi}$, which is (d) at $G=s\phi$, $a=1$. On an arity-zero $O$ the arity-one change along $x_1$ is $-\big\{\{S^{(1)}(x_1),\phi^{(0)}\}+(Q+\{I,-\})\phi^{(1)}(x_1),\,O\big\}$.
\item[(b)] Finite parts of the operators: $O(\varepsilon)\mapsto O(\varepsilon)+\hbar\psi(O)$ for a cochain-valued linear map $\psi$: $L_1\mapsto L_1+[s_{\mathfrak{psl}(4|4)},\psi]$.
\item[(c)] Gauge-fixing operator (proof plan): a smooth family $Q^{GF}_t$ with $[Q,Q^{GF}_t]$ generalised Laplacians commuting with $Q$ and covariant under translations, $\mathfrak{so}(V)$, $\mathfrak{so}(W_{\mathbb R}^\perp)$ is expected to change $L_1$ by $[s_{\mathfrak{psl}(4|4)},\Psi]$ plus changes of types (a), (b), (d).
\item[(d)] Closed finite counterterm $F_{\mathfrak{psl}(4|4)}\mapsto F_{\mathfrak{psl}(4|4)}+\hbar\,a\,G$ with $s_{\mathfrak{psl}(4|4)}G=0$: $L_1\mapsto L_1-a\,\{G,-\}$. The direction of the coupling is
\begin{equation}
G=S-S^{(2)},\qquad \frac{d}{d\mu}\Big|_{\mu=1}\big(\mu S+S^{(1)}+\mu^{-1}S^{(2)}\big),
\label{eq:o-coupling}
\end{equation}
since $\mu S+S^{(1)}+\mu^{-1}S^{(2)}$ solves [ch02, eq:c2-mc] for every $\mu$; it changes $L_1^{(0)}$ by $-a\{S,-\}=-a(Q+\{I,-\})$, which vanishes on gauge-invariant $O$, leaves $L_1^{(1)}$ unchanged and changes $L_1^{(2)}(x_1,x_2)$ by $+a\{S^{(2)}(x_1,x_2),-\}$.
\item[(e)] The outer direction $G_B$ of Proposition \ref{prop:o-GB}, $F_{\mathfrak{psl}(4|4)}\mapsto F_{\mathfrak{psl}(4|4)}+\hbar\,a\,G_B$: $L_1\mapsto L_1-a\{G_B,-\}$; on an arity-zero $O$ the arity-one change along a $B$-eigenvector $x_1$ with eigenvalue $b$ is $-a\,b\,\{S^{(1)}(x_1),O\}$, the arity-two change along $B$-eigenvectors $x_1,x_2$ is $-a(b_1+b_2)\{S^{(2)}(x_1,x_2),O\}$, and $L_1^{(0)}$ is unchanged.
\end{itemize}
On the twisted differential, with $N_{R_\xi}:=\{S^{(1)}(R_\xi),-\}$ for $R_\xi\in\mathfrak{so}(W_{\mathbb R}^\perp)$ with $[R_\xi,\xi]=\xi$ (two-component $r^i{}_j=\tfrac43(w^in_j-\tfrac14\delta^i_j)$ in (B3) of the dictionary):
\begin{equation}
\begin{aligned}
&\text{(a), (b)}:\ Q_1^\xi\mapsto Q_1^\xi+[Q_0^\xi,\Psi^\xi],\qquad \text{(d) with \eqref{eq:o-coupling}}:\ Q_1^\xi\mapsto Q_1^\xi-a\,Q_0^\xi-a\,[N_{R_\xi},Q_0^\xi],\\
&\text{(e)}:\ Q_1^\xi\mapsto Q_1^\xi+a\,[N_{R_\xi},Q_0^\xi],\qquad [N_{R_\xi},Q_0^\xi]=\big\{-S^{(1)}(\xi)+S^{(2)}(\xi,\xi),-\big\}=[N_{\rm gh},Q_0^\xi]-Q_0^\xi ,
\end{aligned}
\label{eq:o-twistlaws}
\end{equation}
with $N_{\rm gh}$ the derivation multiplying a functional of degree $k$ by $k$.
\end{proposition}

\begin{proof}
(a) At order $\hbar$, $O[T]$ changes by minus the tree insertion of $\hbar s_{\mathfrak{psl}(4|4)}\phi$ [ch06, def:c6-tree]. The equivariant form of [ch06, lem:c6-bvexact] (same proof with $d_T$ and Lemma \ref{lem:o-Tflow}) gives, for odd $\phi$,
\[
\mathrm{tree}(s\phi,O)=d_T\,\mathrm{tree}(\phi,O)-\{\phi,O\}[T]+\mathrm{tree}(\phi,sO),\qquad s=s_{\mathfrak{psl}(4|4)} .
\]
Hence $d_TO[T]'=L(O)[T]+\hbar\big(s\{\phi,O\}-\{\phi,sO\}\big)[T]-\hbar\,\mathrm{tree}(s\phi,sO)$, while $L(O)[T]'=L(O)[T]-\hbar\,\mathrm{tree}(s\phi,sO)+O(\hbar^2)$. The derivation rules of the bracket give $s\{\phi,O\}-\{\phi,sO\}=-\{s\phi,O\}$ for odd $\phi$ (checked with the bracket of [ch01, def:c1-bracket] in a finite-dimensional model), i.e.\ $[s,\mathrm{ad}_\phi]=-\mathrm{ad}_{s\phi}$. The arity-one part of $s\phi$ along $x_1$ is $\{S^{(1)}(x_1),\phi^{(0)}\}+(Q+\{I,-\})\phi^{(1)}(x_1)$, the term $d_{CE}\phi^{(0)}$ being zero in arity one. (b) $O[T]'=O[T]+\hbar\psi(O)[T]$, $d_T\psi(O)[T]=(s\psi(O))[T]+O(\hbar)$, and $L(O)[T]'=L(O)[T]+\hbar\psi(sO)[T]$. (d) For even $G$ with $sG=0$, differentiating the flow along $J\mapsto J-\hbar aG$ in the UV-boundary formula gives $d_T\,\mathrm{tree}(G,O)=\{G,O\}[T]+\mathrm{tree}(G,sO)$; the rest is as in (a). The family $\mu S+S^{(1)}+\mu^{-1}S^{(2)}$: by [ch02, lem:c2-components] the arity-$n$ component of the Maurer--Cartan equation for $\sum_kc_kS^{(k)}$ is the original one multiplied by $c_0c_n$ (terms $\{S,S^{(n)}\}$), $c_{n-1}$ (terms $S^{(n-1)}([\cdot,\cdot],\dots)$) and $c_jc_{n-j}$; with $c_0=\mu$, $c_1=1$, $c_2=\mu^{-1}$, $c_{k\ge3}=0$ every product in arities $1,2,3$ is $1$, $\mu$, resp.\ $\mu^{-1}$, uniformly within each arity, and arity $\ge4$ contains $\{S^{(2)},S^{(2)}\}=0$ only. The components of $-a\{G,-\}$: $G^{(0)}=S$ and $\{S,-\}=Q+\{I,-\}$, $G^{(1)}=0$, $G^{(2)}=-S^{(2)}$. (e) Proposition \ref{prop:o-GB}. Twisted forms: $r_\xi(G)|_{q=1}=S+\tfrac12S^{(2)}(\xi,\xi)=S^\xi-S^{(1)}(\xi)+S^{(2)}(\xi,\xi)$; [ch02, lem:c2-components] with $(R_\xi,\xi)$ and $(R_\xi,\xi,\xi)$ gives $\{S^{(1)}(R_\xi),S^{(1)}(\xi)\}=-S^{(1)}(\xi)$, $\{S^{(1)}(R_\xi),S^{(2)}(\xi,\xi)\}=-2S^{(2)}(\xi,\xi)$, $\{S^{(1)}(R_\xi),S\}=0$, hence $N_{R_\xi}S^\xi=-S^{(1)}(\xi)+S^{(2)}(\xi,\xi)$ and $[N_{R_\xi},Q_0^\xi]=\{N_{R_\xi}S^\xi,-\}$, so $\{r_\xi(G),-\}=Q_0^\xi+[N_{R_\xi},Q_0^\xi]$. For (e), $r_\xi(G_B)|_{q=1}=S^{(1)}(\xi)-\tfrac12\cdot2S^{(2)}(\xi,\xi)=-N_{R_\xi}S^\xi$. Finally $N_{\rm gh}S^\xi=-S^{(1)}(\xi)+S^{(2)}(\xi,\xi)=N_{R_\xi}S^\xi$ (degrees $0,-1,-2$), and since the bracket has degree $+1$, $[N_{\rm gh},\mathrm{ad}_f]=\mathrm{ad}_{N_{\rm gh}f}+\mathrm{ad}_f$; with $f=S^\xi$ this is the last identity of \eqref{eq:o-twistlaws}. (c) Proof plan: Duhamel's formula for $\frac{d}{dt}\int_\varepsilon^TQ^{GF}_te^{-sH_t}ds$ splits $\frac{d}{dt}P_t(\varepsilon,T)$ into $(Q\otimes1+1\otimes Q)$ of a kernel and kernels at Schwinger times $\varepsilon$ and $T$; the first is a homotopy of the flow, the second a UV-boundary term whose limit is local and is to be absorbed by (a), (b), (d), the third a change of the endpoint data. Different gauge fixings give canonically homotopy equivalent spaces of quantisations [C07, §1.10]; in this program $Q^{GF}$ is the unique Feynman-gauge operator of [ch03, prop:c3-H].
\end{proof}

\begin{proposition}[The outer direction $G_B$]
\label{prop:o-GB}
Put $G_B:=\frac{d}{dt}\big|_{t=0}\big(e^{tB}\big)^*S_{\mathfrak{psl}(4|4)}$, i.e.
\[
G_B^{(0)}=0,\qquad G_B^{(1)}(x)=S^{(1)}(Bx),\qquad G_B^{(2)}(x_1,x_2)=S^{(2)}(Bx_1,x_2)+S^{(2)}(x_1,Bx_2),\qquad G_B^{(k\ge3)}=0,
\]
two-component: $G_B^{(1)}$ is $S^{(1)}$ with $(\xi_\alpha{}^i,\bar\xi_{\dot\alpha i},\eta^\alpha{}_i,\bar\eta^{\dot\alpha i})$ replaced by $(\xi_\alpha{}^i,-\bar\xi_{\dot\alpha i},-\eta^\alpha{}_i,\bar\eta^{\dot\alpha i})$, zero on even generators, and $G_B^{(2)}$ is $\pm2S^{(2)}$ on pairs from one $B$-eigenspace and $0$ on pairs from different ones; along $\xi=u\otimes w$, $G_B^{(1)}(\xi)$ is \eqref{eq:o-S1xi} and $G_B^{(2)}(\xi,\xi)=-\int(u_\alpha w^i\lambda^{*\alpha}{}_i,u_\beta w^j\lambda^{*\beta}{}_j)$.
\begin{itemize}
\item[(i)] $s_{\mathfrak{psl}(4|4)}G_B=0$, $G_B$ has degree $0$ and weight $4$.
\item[(ii)] $G_B\notin\mathbb C(S-S^{(2)})+\operatorname{Im}s_{\mathfrak{psl}(4|4)}$ on local cochains.
\item[(iii)] Under $F\mapsto F+\hbar aG_B$: $Q_1\mapsto Q_1-a\,\delta_\xi$ on $S_+\otimes W$ and $Q_1\mapsto Q_1+a\,\delta_{\bar\xi}$ on $S_-\otimes W^*$; $S_1\mapsto S_1+a\,\delta_\eta$ on $S_+\otimes W^*$ and $S_1\mapsto S_1-a\,\delta_{\bar\eta}$ on $S_-\otimes W$; $D_1$, $K_1$, $P_1$, $M_1$, $R_1$ and $L_1^{(0)}$ are unchanged; along $(\xi,\xi)$ the arity-two component changes by $-2a\{S^{(2)}(\xi,\xi),-\}$; twisted, $Q_1^\xi\mapsto Q_1^\xi+a[N_{R_\xi},Q_0^\xi]$.
\end{itemize}
\end{proposition}

\begin{proof}
(i) $e^{tB}$ is a family of automorphisms of $\mathfrak{psl}(4|4)$ (Lemma \ref{lem:o-Bouter}); pulling back cochains by an automorphism commutes with $d_{CE}$ and with the bracket of values, so $(e^{tB})^*S_{\mathfrak{psl}(4|4)}$ solves [ch02, eq:c2-mc] for every $t$, and its $t$-derivative at $t=0$ is $s_{\mathfrak{psl}(4|4)}$-closed. $B$ is even and commutes with the dilatation, so the degree and the weight are those of $S_{\mathfrak{psl}(4|4)}$. (ii) Suppose $G_B=aG+s\phi$ with $G$ the coupling direction and $\phi$ local of degree $-1$; take the weight-$4$ part. In arity one along $\xi\in S_+\otimes W$, $G^{(1)}=0$, so $S^{(1)}(\xi)=\{S^{(1)}(\xi),\phi^{(0)}\}+(Q+\{I,-\})\phi^{(1)}(\xi)$ up to the shuffle sign. $\phi^{(1)}(\xi)$ has degree $-2$ and weight $\tfrac92$, and no local functional has these (Lemma \ref{lem:o-actiontrunc}: $m(2)=4$, and no field has weight $\tfrac12$), so $S^{(1)}(\xi)=\pm\{S^{(1)}(\xi),\phi^{(0)}\}$ for all $\xi\in S_+\otimes W$. This equation is equivariant under the compact real form of $\mathrm{Spin}(4)\times SU(4)$ acting on $\xi$ and on $\phi^{(0)}$, so the average of a solution over this group is a solution, and $\phi^{(0)}$ may be taken invariant: by [ch04, lem:c4-enumeration] and the weight count (three legs need weight $\ge5$) it is a combination $\sum_fa_f\int(f,f^*)$ of the five field-type functionals of Proposition \ref{prop:o-arity0}, and $\{S^{(1)}(\xi),\phi^{(0)}\}$ multiplies each term of \eqref{eq:o-S1xi} by the difference of the weights $a_f$ of its field and antifield. The four linear terms $(\Gamma(\xi,\lambda)|_V,A^*|_V)$, $(\Gamma(\xi,\lambda)|_{W_{\mathbb R}^\perp},A^*|_{W_{\mathbb R}^\perp})$, $(\rho(dA|_V)\xi,\lambda^*)$, $(\rho(dA|_{W_{\mathbb R}^\perp})\xi,\lambda^*)$ give the factors $a_{\lambda,S_-\otimes W^*}-a_{A|_V}$, $a_{\lambda,S_+\otimes W}-a_{A|_{W^\perp}}$, $a_{A|_V}-a_{\lambda,S_+\otimes W}$, $a_{A|_{W^\perp}}-a_{\lambda,S_-\otimes W^*}$, which must all equal $\pm1$; their sum is $0$, a contradiction. (iii) Item (e) of Proposition \ref{prop:o-laws} with $b=\pm1$ on the odd summands and $b=0$ on the even part; along $(\xi,\xi)$, $b_1+b_2=2$; twisted, $r_\xi(G_B)|_{q=1}=-N_{R_\xi}S^\xi$.
\end{proof}

\begin{remark}[The supercharge renormalisation and the R-rotation]
\label{rem:o-alpha}
The arity-one finite counterterm $a\,S^{(1)}(\xi)$ along $\xi$ with its arity-two partner $2a\,S^{(2)}(\xi,\xi)$, forced by the $(\xi,\xi)$ component of the master equation, is the restriction to $\xi$ of $G_B$, and also of the exact cocycle $\frac{d}{dt}|_{t=0}(e^{t\,\mathrm{ad}R_\xi})^*S_{\mathfrak{psl}(4|4)}$ (pull-backs by inner automorphisms are $s$-exact). Twisted, both act by $Q_1^\xi\mapsto Q_1^\xi+a[N_{R_\xi},Q_0^\xi]$: at first order this is conjugation by the R-rotation $\alpha$ of [ESW, Def 2.18], under which $S^{(k)}(\xi,\dots,\xi)$ has weight $k$ and $S$ weight $0$, with $N_\alpha=-N_{R_\xi}$; by \eqref{eq:o-twistlaws}, $a[N_\alpha,Q_0^\xi]=aQ_0^\xi-a[N_{\rm gh},Q_0^\xi]$. On untwisted letters $[N_{R_\xi},Q_0^\xi]=-\delta_\xi+\{S^{(2)}(\xi,\xi),-\}$: it vanishes on $c$ and on $A_{ij}w^j$ ($\xi.c=0$, $\xi.(A_{kl}w^l)=0$), and on $(u'\lambda^i)$ modulo $w^i$ it is $-\tfrac14[A^{ik},A_{kj}w^j]$, the image of the term $\tfrac12\varepsilon_{IJK}[\gamma^J,\gamma^K]$ of [BK, (3.50)] up to the sign of the letter map. Both the coupling direction and $G_B$ are of the form of [targets, Lemma lem:inv].
\end{remark}

\begin{remark}[$G_B$ and the conjugation of radial quantisation]
\label{rem:o-GBdagger}
$B$ takes opposite values on $\xi\in S_+\otimes W$ and on the special supercharges $S_+\otimes W^*$ conjugate to it under $X\mapsto X^\dagger$ (directive R5 (a)). Under $F\mapsto F+\hbar aG_B$ the corrections change as $Q_1\mapsto Q_1-a\,\delta_\xi$, $S_1\mapsto S_1+a\,\delta_\eta$ (Proposition \ref{prop:o-GB} (iii)); compatibility with an antilinear conjugation $Q^\dagger=S$ through one loop requires $-\bar a=a$. The real part of $a$ is therefore expected to be fixed by the conjugation test of R5 (c), and the imaginary part is a phase rotation of the supercharges. This is a statement for the conjugation design, to be checked there.
\end{remark}

\begin{proposition}[The arity-zero freedom in the equivariant scheme]
\label{prop:o-arity0}
An arity-zero finite counterterm $G_0$ is admissible in the equivariant scheme exactly when it extends to an $s_{\mathfrak{psl}(4|4)}$-cocycle $G_0+G^{(1)}+G^{(2)}$, i.e.
\[
\{S^{(1)}(x_1),G_0\}=-(Q+\{I,-\})\,G^{(1)}(x_1)\qquad\text{for every }x_1\in\mathfrak{psl}(4|4),
\]
and the higher components. The weight-$4$ directions left free by the non-equivariant equation of Part C (the four $\beta$-directions and the kernel $\operatorname{span}\{H_1-4H_3-2H_6,\,H_2-2H_4\}$ of the cubic solution, six in all) are expected to be exactly the arity-zero parts of the exact field redefinitions $s\,\phi$ with $\phi$ in the span of
\[
\int(A|_V,A^*|_V),\quad \int(A|_{W_{\mathbb R}^\perp},A^*|_{W_{\mathbb R}^\perp}),\quad \int(\lambda^*,\lambda)\big|_{S_+\otimes W},\quad \int(\lambda^*,\lambda)\big|_{S_-\otimes W^*},\quad \int(c^*,c),
\]
two-component $\int(A_m,A^*_m)$, $\tfrac14\int(A^{ij},A^*_{ij})$, $\int(\lambda_\alpha{}^i,\lambda^{*\alpha}{}_i)$, $-\int(\bar\lambda_{\dot\alpha i},\bar\lambda^{*\dot\alpha i})$, $\int(c^*,c)$ [dictionary, prop:pair], and of the coupling direction \eqref{eq:o-coupling}, provided the image in arity zero of the degree-$0$, weight-$4$ cohomology of $s_{\mathfrak{psl}(4|4)}$ on local cochains is $\mathbb C\,S$. This hypothesis is compatible with $G_B$, whose arity-zero part is $0$; hypothesis (H4) of Theorem \ref{thm:o-twist} is the stronger statement that the cohomology itself is spanned by the classes of $S-S^{(2)}$ and $G_B$. Since both sides count six directions, the content of the statement is that the projection to arity zero of the five $s\phi$ and the coupling direction has rank six and image the Part-C directions (E3b). The forced term $\tfrac{27}{16}\,\frac{\hbar\operatorname{tr}_{\mathfrak g}(\mathrm{ad}_X\mathrm{ad}_Y)}{16\pi^2(X,Y)}\int(d^*A|_V,d^*A|_V)$ [ch04, thm:c4-finite] has $\{S^{(1)}(\xi),-\}$- and $\{S^{(1)}(K_{v_1}),-\}$-variations [ch04, prop:c4-K] that enter the arity-one obstructions along $\xi$ and $K_{v_1}$ as $\{S^{(1)}(x_1),F_0\}$.
\end{proposition}

\begin{proof}
The equivariant one-loop master equation changes by $-\hbar\,s_{\mathfrak{psl}(4|4)}G$ under $F\mapsto F+\hbar G$; its arity-one component is the displayed condition. The weight-$4$ degree-$-1$ local functionals with two legs are the five listed (one antifield of weight $\ge\tfrac52$, [ch04, lem:c4-enumeration]); with three legs none exists [Part C, Thm thm:s2c-solve]. The coupling direction is closed by \eqref{eq:o-coupling}, $G_B$ by Proposition \ref{prop:o-GB}. Under the hypothesis, an admissible $G_0$ is the arity-zero part of $cG+s\phi+c'G_B$, i.e.\ $cS+(s\phi)^{(0)}$. The hypothesis and the rank of the projection are checked by the engine (E3b): solve the displayed equation for the Part-C directions with $x_1=\xi$.
\end{proof}

\begin{theorem}[Scheme-independent content]
\label{thm:o-invariants}
Under (a), (b), (d) with the coupling direction, and (e) of Proposition \ref{prop:o-laws}:
\begin{itemize}
\item[(i)] $[Q_1^\xi]$ on $H(\mathscr O_l,Q_0^\xi)$ is invariant; so are its kernel, image, rank and the lifting statements [targets, T8, Lemma lem:inv]. Under (H4) of Theorem \ref{thm:o-twist} the same holds for (d) with every closed $G$.
\item[(ii)] For a gauge-invariant $O$ with $L_1(O)^{(0)}=0$ in the schemes compared, $D_1(O)$ changes by an $s$-exact local functional under (a), (d), (e), and under (b) when $\psi$ preserves the Weyl weight; the coefficient $\gamma_O$ in $D_1(O)\equiv2\gamma_O\,O$ is invariant.
\item[(iii)] $Q_1$ and $S_1$ change jointly: by $-\{(s\phi)^{(1)}(x_1),-\}$ with the same $\phi$ under (a), by $[s,\psi]$ with the same $\psi$ under (b), and by $-a\,b\,\delta_{x_1}$ with opposite signs on $\xi\in S_+\otimes W$ and on $\eta\in S_+\otimes W^*$ under (e). Their pair values are scheme dependent; for instance $c_K$ of Proposition \ref{prop:o-konishi} changes as $c_K\mapsto c_K-\kappa_Bx$ under $\psi(B_w)=x\,\bar F_w$ and is not an entry of $[Q_1^\xi]$, because $Q_0^\xi O_\zeta=\kappa_BB_w\neq0$.
\item[(iv)] Under the hypotheses of Theorem \ref{thm:o-KD}, $K_1(O)\equiv\tfrac14\operatorname{div}v\,\gamma_O\,O$ modulo $s$-exact terms for the operators listed there, in every scheme with $P_1=M_1=R_1=0$.
\end{itemize}
\end{theorem}

\begin{proof}
(i) Theorem \ref{thm:o-twist} (iii). (ii) The arity-one component along $\Delta$ of $-\mathrm{ad}_{s\phi}(O)$ for a cocycle $O$ is $-\{\{S^{(1)}(\Delta),\phi^{(0)}\},O\}-\{(Q+\{I,-\})\phi^{(1)}(\Delta),O\}$; the five functionals of Proposition \ref{prop:o-arity0} have weight $4$, so $\{S^{(1)}(\Delta),\phi^{(0)}\}=0$ (Example c2-dil: $S^{(1)}(\Delta)$ multiplies a density of weight $4$ by $2E+8$, an integrated total derivative), and the second term equals $-(Q+\{I,-\})\{\phi^{(1)}(\Delta),O\}$ on the cocycle $O$, which is $s$-exact. (d) with the coupling direction contributes $-a\{G^{(1)}(\Delta),O\}=0$, $G$ having no arity-one part; (e) contributes $-a\{S^{(1)}(B\Delta),O\}=0$, since $B\Delta=0$. For (b), $[\delta_\Delta,\psi]=0$ when $\psi$ preserves the weight. (iii) Proposition \ref{prop:o-laws} (a), (b), (e) with $x_1=\xi$ and $x_1=\eta$, and Proposition \ref{prop:o-GB} (iii). For $c_K$: $(Q+\{I,-\})O_\zeta=0$ by gauge invariance, $\{S^{(2)}(\xi,\xi),O_\zeta\}$ is proportional to $(u'^\alpha u_\alpha)A_{ij}w^iw^j=0$, and $\{S^{(1)}(\xi),O_\zeta\}=\kappa_BB_w$ (Proposition \ref{prop:o-konishi}). (iv) Theorem \ref{thm:o-KD} with (ii).
\end{proof}

\begin{tabular}{l|l|l}
freedom & $Q_1$ & $S_1$ \\
\hline
$\tfrac{27}{16}$ term (forced) & via $\{S^{(1)}(\xi),F_0\}$ in $F^{(1)}(\xi)$ & via $\{S^{(1)}(\eta),F_0\}$ in $F^{(1)}(\eta)$ \\
field redefinitions (a) & $-\{\{S^{(1)}(\xi),\phi^{(0)}\}+(Q+\{I,-\})\phi^{(1)}(\xi),-\}$ & $-\{\{S^{(1)}(\eta),\phi^{(0)}\}+(Q+\{I,-\})\phi^{(1)}(\eta),-\}$ \\
coupling (d) & arity one: none; arity zero $-a(Q+\{I,-\})$; arity two $+a\{S^{(2)},-\}$ & idem \\
outer derivation (e), $G_B$ & $-a\,\delta_\xi$ on $S_+\otimes W$, $+a\,\delta_{\bar\xi}$ on $S_-\otimes W^*$ & $+a\,\delta_\eta$ on $S_+\otimes W^*$, $-a\,\delta_{\bar\eta}$ on $S_-\otimes W$ \\
arity-one counterterms, $s$-exact & (a) with $\phi^{(0)}=0$ & idem \\
operator finite parts (b) & $+[s,\psi]^{(1)}(\xi)$ & $+[s,\psi]^{(1)}(\eta)$ \\
gauge fixing (c) & proof plan & proof plan \\
\end{tabular}
\tabcaption{Effect of the scheme choices on $Q_1$ and $S_1$ ($F\mapsto F+\hbar aG$); $\delta_x=\{S^{(1)}(x),-\}$ is the classical transformation. Twisted: (a), (b) add $[Q_0^\xi,\Psi^\xi]$, the coupling direction adds $-aQ_0^\xi-a[N_{R_\xi},Q_0^\xi]$, $G_B$ adds $a[N_{R_\xi},Q_0^\xi]$. Classical-type coefficients (E4, the classical-type entries of E5) are shifted by $G_B$ and by the field redefinitions and are reported as families.}

\begin{tabular}{l|l|l|l}
freedom & $D_1$ & $K_1$ & $P_1$, $M_1$, $R_1$ \\
\hline
$\tfrac{27}{16}$ term (forced) & none & via $\{S^{(1)}(K_{v_1}),F_0\}$ [ch04, prop:c4-K] & none ($F_0$ covariant) \\
field redefinitions (a) & $s$-exact & $-\{(Q+\{I,-\})\phi^{(1)}(K_{v_1}),-\}$ for the five $\phi$ & $0$ for the five $\phi$ \\
coupling (d) & none & none & none \\
outer derivation (e), $G_B$ & none & none & none \\
arity-one counterterms, $s$-exact & $s\{\phi^{(1)}(\Delta),O\}$ & $\phi^{(1)}(K_{v_1})\in\mathbb C\int(c^*,\iota_vA)$ & $0$ \\
operator finite parts (b) & $0$ if $\psi$ weight-preserving & $+[s,\psi]^{(1)}(K_{v_1})$ & $0$ for covariant $\psi$ \\
gauge fixing (c) & proof plan & proof plan & proof plan \\
\end{tabular}
\tabcaption{Effect of the scheme choices on $D_1$, $K_1$, $P_1$, $M_1$, $R_1$. For $\xi$ the weight count excludes arity-one degree-$-2$ functionals of weight $\tfrac92$; for even $x_1$ with vector field $v$ the functional $\int(c^*,\iota_vA)$, two-component $\int(c^*,v^mA_m)$, is the only one; for translations, rotations and $\mathfrak{so}(W_{\mathbb R}^\perp)$ there is none.}

\section{Large distance and the transfer endpoint: $S_1$ and $K_1$}

\begin{proposition}[Floating insertions at large distance]
\label{prop:o-ir}
Let a floating insertion of Lemma \ref{lem:o-legs} sit at $y$ and its neighbours at $z_a,z_b$ in a bounded set. With the tails $D(r)=(4\pi^2r^2)^{-1}$ of the propagator at $T=\infty$ [ch03, prop:c3-P], the integrand over $y$ behaves as: $|y|^{-6}$ for $\xi\in\Sigma$ and for translations; $|y|^{-5}$ for $\eta\in\Sigma^*$, for $\Delta$ and for $\mathfrak{so}(V)$; $|y|^{-4}$ for $K_{v_1}$; $\mathfrak{so}(W_{\mathbb R}^\perp)$ has no floating insertion ($v.A|_V=0$). For $K_{v_1}$, with the propagator factors $2$ ($A\otimes A$) and $-1$ ($A^*|_V\otimes c$) removed, the $A|_V\to c$ insertion has
\begin{equation}
\sum_m\big[L_v\big(D(\cdot-z_a)\,e_k\big)\big]_m(y)\,\partial_{y_m}D(y-z_b)=\frac{(b\cdot y)\,y_k-|y|^2b_k}{4\pi^4|y|^6}+O(|y|^{-5}),\qquad b=v_1,
\label{eq:o-Kir}
\end{equation}
with average over the sphere $|y|=R$ equal to $-\dfrac{3}{16\pi^4}\dfrac{b_k}{R^4}\neq0$. Two-component: the left side is $\sum_m\big(v^n\partial_n(\delta_{mk}D_a)+D_a\partial_mv^k\big)\partial_mD_b$ with $v^m=y^2b^m-2(b\cdot y)y^m$. A $P$-edge pair of floating odd insertions at $y_1,y_2$, with spinor data of degrees $g_1,g_2$ in the position ($g=0$ for $\Sigma$, $g=1$ for $\Sigma^*$), has, for one orientation, superficial degree $g_1+g_2-1$ at large $(y_1,y_2)\in\mathbb R^8$. At finite $T$ every component of $P(\varepsilon,T)$ and its derivatives are bounded by the corresponding $T=\infty$ tails, uniformly in $T$, and decay as $e^{-|y|^2/4T}$ for $|y|^2\gg T$.
\end{proposition}

\begin{proof}
Coefficient growth: $|\epsilon(y)|\sim1$, $|y|$; $|v(y)|\sim|y|$ for $\Delta$, $|y|^2$ for $K_{v_1}$, with $|dv|\sim1$, $|y|$. Line decay: $D\sim|y|^{-2}$, $\partial D,{\not d}D\sim|y|^{-3}$ (Lemma \ref{lem:o-legs}). Products: $1\cdot|y|^{-3}\cdot|y|^{-3}$; $|y|\cdot|y|^{-3}\cdot|y|^{-3}$; $|y|^2\cdot|y|^{-3}\cdot|y|^{-3}$ and $|y|\cdot|y|^{-2}\cdot|y|^{-3}$. For \eqref{eq:o-Kir} at $z_a=z_b=0$: $v\cdot\partial D=\tfrac{2(b\cdot y)}{4\pi^2|y|^2}$, $\partial_kD=-\tfrac{2y_k}{4\pi^2|y|^4}$, $y_m\partial_mv^k=2v^k$, and the sum is $\frac{4}{(4\pi^2)^2}\big((b\cdot y)y_k-|y|^2b_k\big)|y|^{-6}$; the sphere average of $y_ky_l$ is $\tfrac14R^2\delta_{kl}$. The terms linear in $z_a,z_b$ are of order $|y|^{-5}$. An exact symbolic evaluation with exact rational anchors $z_a,z_b$ (three choices) gives the sphere average $-3(4\pi^2)^{-2}b_k$ exactly, and a momentum-space evaluation of the same insertion gives the same coefficient. $P$-edge pair: the chain $c(z)$--$A^*|_V$--$y_1$--$\lambda\lambda$--$y_2$--$A^*|_V$--$c(z')$ has three propagators $\partial D$, ${\not d}D$, $\partial D$ of order $|y|^{-3}$ each and the coefficients $|y_1|^{g_1}|y_2|^{g_2}$; with the measure $d^8y$ the degree is $8-9+g_1+g_2$. Finite $T$: $\partial^k_x\int_\varepsilon^T(4\pi t)^{-2}e^{-|x|^2/4t}dt$ is bounded by $\int_0^\infty|\partial^k_x(4\pi t)^{-2}e^{-|x|^2/4t}|\,dt\le C_k|x|^{-2-k}$ for every $T$, and for $|x|^2\gg T$ the integrand is bounded by $t^{-2-k/2}$ times $e^{-|x|^2/4T}$ on $t\le T$.
\end{proof}

\begin{theorem}[Large-distance behaviour of the graphs]
\label{thm:o-ir}
\begin{itemize}
\item[(i)] For $x_1\in\Sigma\oplus\Sigma^*$, $\Delta$, translations, $\mathfrak{so}(V)$, $\mathfrak{so}(W_{\mathbb R}^\perp)$, every graph of Theorem \ref{thm:o-classes} converges at $T=\infty$ at generic leg momenta, and the graph sums at scale $T$ converge to those at $T=\infty$ as $T\to\infty$.
\item[(ii)] For $x_1=K_{v_1}$, classes (T), (A), (B), (E1), (E2) converge at $T=\infty$. A graph of class (C) or (D) with $S^{(1)}(K_{v_1})$ floating on the cycle has, at finite $T$, a term $\log T$ whose coefficient is $3\cdot2\pi^2(4\pi^2)^{-2}=\tfrac{3}{8\pi^2}$ times the tree obtained by cutting the insertion out of the cycle and evaluating the freed legs on the constant modes $b\otimes(-,-)^{-1}$ (end $A|_V$) and $1$ (end $c$) (half the coefficient of $\log R$ of a radial cutoff $|y|<R$); this tree contains the special edge at a generic momentum $p$ and is $O(\varepsilon)$, through $e^{-\varepsilon|p|^2}-1$. The graph diverges as $T\to\infty$ at fixed $\varepsilon>0$ when the tree is nonzero; the nonvanishing of the tree is not established.
\item[(iii)] The arity-one parts $O[T]^{(1)}(K_{v_1})$ and $W(P(\varepsilon,T),I_{\mathfrak{psl}(4|4)})^{(1)}(K_{v_1})$ contain cycles with $S^{(1)}(K_{v_1})$ floating and no special edge; their $\log T$ coefficient is the two-zero-mode tree $Y$ of the same construction, of order $\varepsilon^0$. Both objects exist at every finite $T$, and their limits $T\to\infty$ exist exactly when $Y=0$ in the corresponding sector.
\item[(iv)] Arity two. A $P$-edge pair of floating $S^{(1)}(\eta_1)$, $S^{(1)}(\eta_2)$, $\eta_a\in\Sigma^*$, has after the sum over the two orientations and the two chirality assignments a logarithmic large-distance term ($\log T$ at finite $T$): the orientation sum cancels the degree-$1$ part identically, and the surviving degree-$0$ part is proportional to $v_1=\Gamma(\eta_1,\eta_2)|_V$, two-component $v_1^m=-i\big(\eta_2{}^\alpha{}_i\sigma^m_{\alpha\dot\alpha}\bar\eta_1^{\dot\alpha i}+\eta_1{}^\alpha{}_i\sigma^m_{\alpha\dot\alpha}\bar\eta_2^{\dot\alpha i}\big)$, times the first moment of the rest of the cycle: in momentum space, with the fibre factors removed and the rest attached through $\hat f(q)=f_0+f_1\cdot q+\cdots$ at the first end and $\hat g(p)=g_0+g_1\cdot p+\cdots$ at the second, the $S^3$ average of the degree-$(-4)$ symbol is $6i\,v_1\cdot(f_0g_1-f_1g_0)$. The coefficient is $O(\varepsilon)$ in defect graphs (special edge in the rest) and of order $\varepsilon^0$ in $O[T]^{(2)}(\eta_1,\eta_2)$ and $W(P(\varepsilon,T),I_{\mathfrak{psl}(4|4)})^{(2)}(\eta_1,\eta_2)$; it vanishes, pointwise in the degree-$0$ symbol, exactly when $v_1=0$. $P$-edge pairs $(\xi,\xi')$ (degree $-1$) and $(\xi,\eta)$ (degree $0$ for one orientation, lowered to $-1$ by the orientation sum) converge, and so does every other configuration of the arity-two graphs along $(\xi,\xi')$, $(\xi,\eta)$ and along $(\eta_1,\eta_2)$ with $v_1=0$.
\end{itemize}
\end{theorem}

\begin{proof}
(i) Proposition \ref{prop:o-ir}; vertices with external legs or attached to a special edge or to $O$ are anchored, because every vertex of $I$ carries a leg or a tree to the legs and $K_\varepsilon-K_0$ has range $\sqrt\varepsilon$; at arity one there is a single floating insertion. The limit $T\to\infty$ follows by dominated convergence with the uniform bounds of Proposition \ref{prop:o-ir}. (ii) In (T), (A), (B) the vertex $S^{(1)}(K_{v_1})$ is merged with $O$ or is an end of a special edge; in (E1), (E2) the graphs are trees, and a bridge carries a generic momentum $p\neq0$, where $|p|^{-4}$ is finite. In (C), (D) the floating vertex is bivalent on the cycle; the tail \eqref{eq:o-Kir} has nonzero sphere average and is independent of $z_a,z_b$ at leading order, so the $y$-integral equals the stated multiple of the logarithm times the graph with the constant kernel $b\otimes(-,-)^{-1}$ in place of the composite line (factors $2\cdot(-1)\cdot(-3)=6$, $\int_{1<|y|<R}|y|^{-4}d^4y=2\pi^2\log R$, and $\int|y|^{-4}\varphi(|y|^2/T)\,d^4y=\pi^2\log T+O(1)$ for the Gaussian cutoff profile $\varphi$ of $P(\varepsilon,T)$, $\varphi(0)=1$). Cutting the cycle there leaves a tree in which every edge carries a sum of leg momenta of a nonempty proper subset, generic, and the special edge contributes $e^{-\varepsilon|p|^2}-1=O(\varepsilon)$. (iii) Same computation without special edge; at finite $T$ the cutoff profile makes every integral finite. (iv) Two-component, with $\lambda$ at $y_1$ and $\bar\lambda$ at $y_2$, $\bar\epsilon_1(y)=i\,y_a\eta_1\sigma^a$, $\epsilon_2(y)=-i\,y_b\sigma^b\bar\eta_2$ (Definition \ref{def:o-names}), the chain is, up to the fibre factors,
\begin{equation}
C(\eta_1,\eta_2)[z,z']=\int\!\!\int\partial_mD(y_1-z)\,\big(\bar\epsilon_1(y_1)\bar\sigma^m\sigma^k\bar\sigma^n\epsilon_2(y_2)\big)\,\partial_kD(y_1-y_2)\,\partial_nD(y_2-z')\,d^4y_1\,d^4y_2 .
\label{eq:o-pair}
\end{equation}
The two contractions of the $A^*|_V$ legs with the two ghost ends differ by the exchange of the odd ends (Koszul sign $-1$, [ch03, prop:c3-P]), and the other chirality assignment is the relabelling $\eta_1\leftrightarrow\eta_2$, so the total is $C(\eta_1,\eta_2)[z,z']-C(\eta_1,\eta_2)[z',z]+(1\leftrightarrow2)$; the symmetry under $1\leftrightarrow2$ is the graded symmetry of the arity-two cochain in its odd arguments. In momentum space the single-orientation degree-$(-5)$ symbol is $4\,(\eta_1{\not p}\bar\eta_2)/|p|^6$, odd, and cancels in the difference; the difference is $(f_0g_1-f_1g_0)_c$ times a degree-$(-4)$ symbol, pointwise $\big(8(\eta_1{\not p}\bar\eta_2)p_c|p|^{-2}+4(\eta_1\sigma^c\bar\eta_2)\big)|p|^{-4}$, with $S^3$ average $6(\eta_1\sigma^c\bar\eta_2)$, and the sum over $1\leftrightarrow2$ gives $6i\,v_1^c$. Shell integrals over $\rho<|p|<2\rho$ with an exact Gaussian rest converge to $2\pi^2\log2$ times $6i\,v_1\cdot(f_0g_1-f_1g_0)$ as $\rho\to0$. For $v_1=0$ the degree-$(-4)$ symbol vanishes pointwise. For $(\xi,\eta)$ the two orientations have equal IR-leading symbols with the $\lambda\lambda$ edge traversed in opposite directions, so the difference vanishes identically, while one orientation alone has the nonzero $S^3$ average $2\,\eta^\alpha{}_i\xi_\alpha{}^i$ (fibre factors removed). These statements are machine-checked numerically on the explicit Euclidean $\sigma$ matrices. Special-edge pairs: $\int du\,(K_\varepsilon-K_0)(u)\,g(y+u,y)=\big((e^{\varepsilon\Delta_u}-1)g(y+u,y)\big)_{u=0}=O(\varepsilon|y|^{-6})$; $S^{(2)}$ is never floating (Lemma \ref{lem:o-legs}).
\end{proof}

\begin{definition}[Components along $K$]
\label{def:o-K}
$K_1(O):=\hbar L_1(O)^{(1)}(K_{v_1})$, where $L_1(O)^{(1)}(K_{v_1})$ is the local representative at a finite transfer endpoint $T$: the $\varepsilon\to0$ limit at fixed $T$ of the K-route graphs of Theorem \ref{thm:o-classes} with $x_1=K_{v_1}$ and edges $P(\varepsilon,T)$, read off by \eqref{eq:o-L1graphs}. The arity-two components along $(\eta_1,\eta_2)$ and $(K_{v_1},P_{v_2})$ are defined in the same way at finite $T$.
\end{definition}

\begin{proposition}[Placeholder: existence and $T$-independence along $K$]
\label{prop:o-Kclaim}
Claim, to be proved or refuted (not established in this design). For every $T\in(0,\infty)$ and every operator of the target table, the $\varepsilon\to0$ limit at fixed $T$ of the K-route graphs exists and equals $\hbar L_1(O)^{(1)}(K_{v_1})[T]_{\rm tree}$ plus the classical and arity-zero terms of Theorem \ref{thm:o-classes}, with $L_1(O)^{(1)}(K_{v_1})$ a local functional independent of $T$. The same for the arity-two components along $(\eta_1,\eta_2)$ and along $(K_{v_1},P_{v_2})$.
\end{proposition}

\begin{remark}[What the claim requires]
By Proposition \ref{prop:o-Tindep}, existence of a local representative at one finite $T$, together with the master equation, gives the same local representative at every finite $T$; the content of the claim is therefore the existence at fixed $T$. The expected mechanism (directive R7): the classes (T), (A), (B), (E1), (E2) converge at $T=\infty$ (Theorem \ref{thm:o-ir} (ii)); the large-distance terms of the classes (C), (D) are $\log T$ times $O(\varepsilon)$ coefficients, which vanish as $\varepsilon\to0$ at fixed $T$; the surface terms of the classical master equation vanish at finite $T$ (Section 3); the singular parts are independent of $T$ (Lemma \ref{lem:o-Tflow} (iv)), and by [ch04, lem:c4-weights] extended with the weights $-1,0,+1$ of $v$, $dv$, $\partial^2v$ the coefficient of $\varepsilon^{-1}$ along $K_{v_1}$ vanishes and that of $\log\varepsilon$ lies in the span of $(v.A,A^*)$, $(v.\lambda,\lambda^*)$, $(v.c,c^*)$-type structures of weight $4$. What remains to be shown is the locality of the $\varepsilon\to0$ limit at fixed $T$ and the uniformity in $T$ of the $O(\varepsilon\log T)$ bound. Once the claim holds, $K_1$ is the same for every $T$ and the limit $T\to\infty$ of $K_1$ is trivially safe; the observables $O[T]^{(1)}(K_{v_1})$ themselves keep the $\log T$ term of Theorem \ref{thm:o-ir} (iii) when $Y\neq0$, as objects at the transfer endpoint.
\end{remark}

\begin{proposition}[Cross-checks of $K_1$]
\label{prop:o-Kconsistent}
Assume Proposition \ref{prop:o-Kclaim}. Write $C_T(O):=d_TO[T]$.
\begin{itemize}
\item[(i)] For $\eta_1,\eta_2\in\Sigma^*$ with $[\eta_1,\eta_2]=-K_{v_1}$, $v_1=\Gamma(\eta_1,\eta_2)|_V$ [ch02, eq:c2-oddodd], the arity-two component of \eqref{eq:o-cocycle} at $(\eta_1,\eta_2)$ contains $L_1(O)^{(1)}(K_{v_1})$ through its $d_{CE}$ term $(-1)^{|O|+1}L_1(O)^{(1)}([\eta_1,\eta_2])=(-1)^{|O|}\hbar^{-1}K_1(O)$ and expresses $K_1(O)$ through $S_1$, the arity-two component $L_1(O)^{(2)}(\eta_1,\eta_2)$, $L_1(O)^{(0)}$ and $L_1(s_{\mathfrak{psl}(4|4)}O)$; the remaining terms depend on $(\eta_1,\eta_2)$ through $v_1$ only.
\item[(ii)] At $T=\infty$, for pairs with $v_1=0$, $O[\infty]^{(1)}(\eta_a)$ and $O[\infty]^{(2)}(\eta_1,\eta_2)$ converge (Theorem \ref{thm:o-ir} (i), (iv): the degree-$0$ symbol of the $P$-edge pair vanishes pointwise when $v_1=0$), and $(d_{CE}+Q)C_\infty(O)=0$ in arity two reads $QC_\infty(O)^{(2)}(\eta_1,\eta_2)=0$. The products $\eta_1\eta_2$ with $\Gamma(\eta_1,\eta_2)|_V=0$ span the kernel of $\Gamma|_V$ on $\mathrm{Sym}^2\Sigma^*$ (rank $132=136-4$ in the realisation of Chapter 2).
\item[(iii)] With the one-loop inner product of directive R5 and $P_1=0$, the relation $K=P^\dagger$ gives a second determination of $K_1$; it is a cross-check, and it becomes the definition only if Proposition \ref{prop:o-Kclaim} fails.
\end{itemize}
\end{proposition}

\begin{proof}
(i) The arity-two component of \eqref{eq:o-cocycle} at $(\eta_1,\eta_2)$, with the $d_{CE}$ term of Theorem \ref{thm:o-cocycle} and $[\eta_1,\eta_2]=-K_{v_1}$; the left side depends on $(\eta_1,\eta_2)$ through $[\eta_1,\eta_2]$, so the identity forces the same dependence on the right side. (ii) At $T=\infty$, $d_\infty=d_{CE}+Q$, and $(d_{CE}+Q)C_\infty(O)=0$ in arity two at $(x_1,x_2)$ reads $(-1)^{|C|}C_\infty^{(1)}([x_1,x_2])+QC_\infty^{(2)}(x_1,x_2)=0$ [ch02, eq:c2-dce], $|C|=|O|+1$; for $v_1=0$ the bracket vanishes. Convergence: Theorem \ref{thm:o-ir} (iv) for the $P$-edge pairs, (i) for single floating $S^{(1)}(\eta_a)$. The rank is a machine check. (iii) Directive R5; the inner product is designed separately.
\end{proof}

\begin{proposition}[Where $T\to\infty$ is safe]
\label{prop:o-Tsafe}
\begin{itemize}
\item[(i)] Safe (the graph sums at scale $T$ converge to those at $T=\infty$; the local representatives are those at $T=\infty$ by Proposition \ref{prop:o-Tindep}): arity zero; arity one along $\Sigma$, $\Sigma^*$, $\Delta$, translations, $\mathfrak{so}(V)$, $\mathfrak{so}(W_{\mathbb R}^\perp)$; arity two along $(\xi,\xi')$, $(\xi,\eta)$ and along $(\eta_1,\eta_2)$ with $\Gamma(\eta_1,\eta_2)|_V=0$. In particular $Q_1$, the twisted $Q_1^\xi$ (Theorem \ref{thm:o-twist}), $S_1$ (Proposition \ref{prop:o-S}) and $D_1$ are computed at $T=\infty$.
\item[(ii)] Not established: arity one along $K_{v_1}$ (classes (C), (D) at fixed $\varepsilon$; $O[T]^{(1)}(K_{v_1})$ and $W^{(1)}(K_{v_1})$ when $Y\neq0$); arity two along $(\eta_1,\eta_2)$ with $v_1\neq0$; arity two along $(K_{v_1},P_{v_2})$. These components are defined at finite $T$ (Definition \ref{def:o-K}).
\item[(iii)] Along $\Delta$ the scaling identity of Part D at finite $T$ reads $(M\otimes1+1\otimes M)P(\varepsilon,T)=-4\big(\varepsilon\partial_\varepsilon+T\partial_T\big)P(\varepsilon,T)=4\varepsilon L_\varepsilon-4TL_T$, with the endpoint term $-4TL_T$; by (i) the scaling route of Part D is used at $T=\infty$, where this term is absent.
\end{itemize}
\end{proposition}

\begin{proof}
(i) Theorem \ref{thm:o-ir} (i), (iv) and dominated convergence with the uniform bounds of Proposition \ref{prop:o-ir}; for the $P$-edge pairs the orientation sum is taken before the bound. (ii) Theorem \ref{thm:o-ir} (ii)--(iv). (iii) $\partial_\varepsilon P(\varepsilon,T)=-L_\varepsilon$, $\partial_TP(\varepsilon,T)=L_T$ with $L_t=(Q^{GF}\otimes1)K_t$, and the scaling covariance $\sigma_\lambda P(\varepsilon,T)=P(\lambda^2\varepsilon,\lambda^2T)$ of DEFINITIONS item 6 extended to finite $T$.
\end{proof}

\begin{theorem}[$K_1$ from $D_1$]
\label{thm:o-KD}
Assume $P_1=M_1=R_1=0$ (covariant choice, [ch07, prop:c7-cov]), the cochain identity \eqref{eq:o-cocycle} at $(K_{v_1},P_{v_2})$, and the existence of local representatives along $K_{v_1}$ and at $(K_{v_1},P_{v_2})$ (Proposition \ref{prop:o-Kclaim}). Let $O$ be a scalar bilinear (gauge invariant, $L_1(O)^{(0)}=0$) or a letter; for the letters, read the statements modulo terms built from $(Q+\{I,-\})O$ [ch06, rem:c6-AV]. If $D_1(O)\equiv2\gamma_O\,O$ modulo $s$-exact terms, then
\begin{equation}
K_1(O)\equiv\tfrac14\operatorname{div}v\;\gamma_O\,O\qquad\text{modulo }s\text{-exact terms},
\label{eq:o-KD}
\end{equation}
two-component $K_1(O)\equiv-2\,(b\cdot x)\,\gamma_O\,O$. With the coefficients of [ch06, thm:c6-dil] (the identification of Part D):
\[
\begin{aligned}
&\tfrac{1}{16}\int\omega\,(w_1^{ij}A_{ij},w_2^{kl}A_{kl})-\tfrac1{96}w_1^{ij}w_{2ij}\int\omega\,(A^{kl},A_{kl}):\quad K_1=0;\\
&\tfrac14\int\omega\,(A^{ij},A_{ij}):\quad K_1\equiv48\,(b\cdot x)\,\tfrac{\hbar\operatorname{tr}_{\mathfrak g}(\mathrm{ad}_X\mathrm{ad}_Y)}{16\pi^2(X,Y)}\,O;\\
&\int\omega\,(X,A_{ij}),\ \int\omega\,(X,A_m):\quad 8\,(b\cdot x)\,\tfrac{\hbar\operatorname{tr}_{\mathfrak g}(\mathrm{ad}_X\mathrm{ad}_Y)}{16\pi^2(X,Y)}\,O;\\
&\int\omega\,(\zeta,\lambda):\quad 16\,(b\cdot x)\,\tfrac{\hbar\operatorname{tr}_{\mathfrak g}(\mathrm{ad}_X\mathrm{ad}_Y)}{16\pi^2(X,Y)}\,O;\qquad \int\omega\,(X,c):\quad 4\,(b\cdot x)\,\tfrac{\hbar\operatorname{tr}_{\mathfrak g}(\mathrm{ad}_X\mathrm{ad}_Y)}{16\pi^2(X,Y)}\,O,
\end{aligned}
\]
in the unit of [ch04]. This is a conditional algebraic consequence of the hypotheses.
\end{theorem}

\begin{proof}
The arity-two component of \eqref{eq:o-cocycle} at $(K_{v_1},P_{v_2})$, with $L_1^{(1)}(P_{v_2})=0$ and $L_1^{(1)}$ vanishing on rotations, reads $\pm(v_1,v_2)D_1(O)\pm[\delta_{P_{v_2}},K_1](O)=s(\cdots)$, since $[K_{v_1},P_{v_2}]=(v_1,v_2)\Delta+(\text{rotation})$ [ch07, rem:c7-KDelta] and $S^{(2)}$ vanishes on even arguments; the right side consists of $s$-exact local functionals, which may depend on $x$. By [ch07, cor:c7-K], $K_1(O)=\int\omega\big(v(x)\cdot V+\tfrac14\operatorname{div}v(x)\,\kappa+(\text{antisymmetric part of }dv(x))\cdot R+(\partial^2v)\cdot T_2\big)$ with densities $V,\kappa,R,T_2$ of weights $w_O+1,w_O,w_O,w_O-1$, and $T_2=0$ for these operators. Writing $\delta_{P_{v_2}}$ on $\int\omega$ through integration by parts, $[\delta_{P_{v_2}},K_1](O)=-\int\omega\big((dv(x)v_2)\cdot V+\tfrac14(\partial_{v_2}\operatorname{div}v)\,\kappa+(\partial_{v_2}dv)\cdot R\big)$, with $\tfrac14\partial_{v_2}\operatorname{div}v=-2(v_1,v_2)$. Comparing the dependence on $x$ modulo $s$-exact local functionals, the term with $V$ is $s$-exact, $V\equiv0$ modulo $s$-exact terms; the $R$ term is the rotation component, zero; hence $\kappa$ is a fixed multiple of $\gamma_O\,O$ modulo $s$-exact terms. The multiple is that of the classical relation: $v\mapsto\tfrac14\operatorname{div}v$ is a $1$-cocycle of the conformal Killing fields with values in functions, $\tfrac14\operatorname{div}[v,v']=v(\tfrac14\operatorname{div}v')-v'(\tfrac14\operatorname{div}v)$, so for every number $\gamma$ the assignment $v\mapsto\gamma\,\tfrac14\operatorname{div}v$ satisfies the same relation as the classical Weyl-weight term $w_O\tfrac14\operatorname{div}v$ of $S^{(1)}(v)$ [ch07, lem:c7-weights]; at $v=2E$, $\tfrac14\operatorname{div}(2E)=2$ gives $D_1=2\gamma_O\,O$, hence \eqref{eq:o-KD}. The values: $\tfrac14\operatorname{div}v=-2(b\cdot x)$ and the coefficients $0,-24,-4,-4,-8,-2$ of [ch06, thm:c6-dil]. The cocycle property, $[K_{v_1},P_{v_2}]$ and the values are machine-checked.
\end{proof}

\begin{proposition}[Special supercharges]
\label{prop:o-S}
\begin{itemize}
\item[(i)] Every graph of the arity-one component along $\eta\in\Sigma^*$ converges at $T=\infty$, and $T\to\infty$ is safe (Theorem \ref{thm:o-ir} (i), Proposition \ref{prop:o-Tsafe}); the vertices are [dictionary, (F2)]
\[
S^{(1)}(\eta)=S^{(1)}(\xi)\big|_{\xi\to\epsilon(x)}+\int\Big\{\big(A^{kj}\varepsilon_{\alpha\beta}\eta^\beta{}_j,\lambda^{*\alpha}{}_k\big)-\big(A_{kj}\varepsilon_{\dot\beta\dot\alpha}\bar\eta^{\dot\alpha j},\bar\lambda^{*\dot\beta k}\big)\Big\},
\]
index-free [ch02, eq:c2-S1S], with $\epsilon(x)$ of Definition \ref{def:o-names} at $\xi=0$.
\item[(ii)] With $P_1=0$ and \eqref{eq:o-cocycle} at $(P_{v_2},\eta)$, $[P_{v_2},\eta]=\rho(v_2)\eta\in\Sigma$ [ch02, prop:c2-special]:
\begin{equation}
S_1(\eta)(O)\equiv Q_1\big(\rho(x)\eta\big)(O)+U_\eta(O)\qquad\text{modulo }s\text{-exact terms},
\label{eq:o-SQ}
\end{equation}
where $Q_1(\xi')$ is the one-loop component along $\xi'\in\Sigma$ evaluated pointwise at the Killing spinor, and $U_\eta(O)$ is translation invariant, linear in ${\not d}(\rho(x)\eta)=4\eta$, of weight $w_O-\tfrac12$. $U_\eta=0$ for the scalar bilinears [ch07, cor:c7-S]; for $B_w$, $\bar F_w$ of weight $3$ it lies in the span of the weight-$\tfrac52$ singlets $\int\omega\,(\lambda,\rho(A|_{W_{\mathbb R}^\perp})\zeta)$, $\zeta\in\Sigma$ built from $\eta$, $w$, $u'$, two-component [dictionary, (R3), (P1)]
\[
\int\omega\,(\lambda,\rho(A|_{W_{\mathbb R}^\perp})\zeta)=\int\omega\Big[\big(\lambda_\alpha{}^j,A_{ji}\big)\zeta^{\alpha i}-\big(\bar\lambda_{\dot\alpha j},A^{ji}\big)\bar\zeta^{\dot\alpha}{}_i\Big].
\]
\item[(iii)] For $\eta$ in the commutant of $\xi=u\otimes w$, i.e.\ $\eta^\alpha{}_i=u^\alpha m_i$ with $m_iw^i=0$, or $\eta\in S_-\otimes W$ [targets, T10], the arity-two component of \eqref{eq:o-cocycle} at $(\xi,\eta)$ restricted by $r_\xi$ gives $\{Q_0^\xi,S_1^\xi(\eta)\}+\{S_0^\xi(\eta),Q_1^\xi\}\equiv0$ on $H(Q_0^\xi)$: the BV form of [BK, footnote 24], with $S_1^\xi(\eta)$ the one-loop $R_1$ of the $\mathfrak{psl}(3|3)$ generators, which [BK] do not write out.
\end{itemize}
\end{proposition}

\begin{proof}
(ii) As in the proof of Theorem \ref{thm:o-KD}: by [ch07, prop:c7-select] $S_1(\eta)(O)=\int\omega\big(\mathcal W(\rho(x)\eta)+U_\eta\big)$ with $\mathcal W$ linear in the value of the Killing spinor; $[\delta_{P_{v_2}},S_1(\eta)](O)=-\int\omega\,\mathcal W(\rho(v_2)\eta)$, which the cochain identity equates with $Q_1(\rho(v_2)\eta)(O)$ up to the sign of the classical relation, fixed by the classical components ($S^{(1)}(\eta)$ is $S^{(1)}(\xi)$ at $\xi\to\rho(x)\eta$ plus the term in $4\eta$, [ch02, eq:c2-S1S]); $\rho(v_2)\eta$ spans $\Sigma$. (iii) $[\xi,\eta]=0$ for these $\eta$ by (B1), (B2) of the dictionary [targets, T10 (c)]; apply $r_\xi$ to the arity-two component.
\end{proof}

\section{The Konishi target and the engine plan}

\begin{proposition}[The Konishi entry]
\label{prop:o-konishi}
Let $\zeta,\xi\in S_+\otimes W$, $O_\zeta:=\int\omega\,(\lambda,\rho(A|_{W_{\mathbb R}^\perp})\zeta)$. The arity-one component along $\xi$ at order $\hbar$ has the form
\begin{equation}
\begin{aligned}
Q_1(O_\zeta)={}&c_B'\,\frac{\hbar\operatorname{tr}_{\mathfrak g}(\mathrm{ad}_X\mathrm{ad}_Y)}{16\pi^2(X,Y)}\,B(\xi,\zeta)+c_K\,\frac{\hbar\operatorname{tr}_{\mathfrak g}(\mathrm{ad}_X\mathrm{ad}_Y)}{16\pi^2(X,Y)}\,\tfrac12\int\omega\,\big(\Gamma(\xi,\lambda)|_V,\Gamma(\zeta,\lambda)|_V\big)\\
&+(\text{terms vanishing for }\xi=u\otimes w,\ \zeta=u'\otimes w),
\end{aligned}
\label{eq:o-konishi}
\end{equation}
with $B(\xi,\zeta)=-2\int\omega\,(\rho(\tfrac12[A|_{W_{\mathbb R}^\perp}\wedge A|_{W_{\mathbb R}^\perp}])\xi,\rho(A|_{W_{\mathbb R}^\perp})\zeta)$. For $\xi=u\otimes w$, $\zeta=u'\otimes w$:
\begin{equation}
O_\zeta=\int\omega\,(u'^\alpha\lambda_\alpha{}^i,A_{ij}w^j),\qquad B_w=\int\omega\,\big([A^{ik},A_{kj}w^j],A_{il}w^l\big),\qquad \bar F_w=\int\omega\,\big(\bar\lambda_{\dot\alpha i}w^i,\bar\lambda^{\dot\alpha}{}_jw^j\big),
\label{eq:o-BF}
\end{equation}
$B(\xi,\zeta)=B_w$, $\tfrac12\int\omega\,(\Gamma(\xi,\lambda)|_V,\Gamma(\zeta,\lambda)|_V)=\bar F_w$, and the classical component is $\delta_\xi O_\zeta=\kappa_B\,B_w$ with $\kappa_B=\pm\tfrac14$, the sign fixed by the bracket of [ch01, def:c1-bracket] (engine E0). These are the only singlets of weight $3$ and degree $0$ proportional to $w\otimes w$ and to $u'^\alpha u_\alpha$. The number $c_K$ is the untwisted Konishi entry [targets, T5]: a pair value, with $Q_0^\xi O_\zeta=\kappa_BB_w\neq0$, which changes as $c_K\mapsto c_K-\kappa_Bx$ under the finite part $\psi(B_w)=x\,\bar F_w$ of the weight-$3$ operator $B_w$; no normalisation of the weight-$3$ operators is singled out by the program, and an orthogonality condition to the protected combination leaves the freedom $B_w\mapsto B_w+y\,(\text{protected combination})$ open.
\end{proposition}

\begin{proof}
Selection rules [ch07, prop:c7-select]: weight $w_O+\tfrac12=3$, degree $0$, singlet, linear in $\xi$ and in $\zeta$. Field contents of weight $3$: $AAA$, $\lambda\lambda$ (two dotted or two undotted), $F_AA$, $A\,\partial A$, $c\,A^*$. The $W$-structure $w^iw^j$ is symmetric; $A_{ij}w^iw^j=0$ removes $F_AA$, $A\,\partial A$ and $c\,A^*_{ij}w^iw^j$; undotted $\lambda^i\lambda^j$ carry upper $W$ indices and cannot contract $w\otimes w$; this leaves $B_w$ and $\bar F_w$, whose spinor factor is $u_\alpha u'^\alpha=u'^\alpha u_\alpha=1$. Two-component identities: (R3), (P1) of the dictionary give $(\rho(A)\zeta)^\alpha{}_j=A_{ji}u'^\alpha w^i$, $(\rho(\tfrac12[A\wedge A])\xi)_\alpha{}^i=-\tfrac12[A^{ik},A_{kj}]u_\alpha w^j$ and the pairing $-\tfrac12(u_\alpha u'^\alpha)B_w$; with $\Gamma^m(\xi,\lambda)=i\bar\lambda_{\dot\alpha i}w^i\bar\sigma^{m\dot\alpha\alpha}u_\alpha$ (G1) and the Euclidean identity $\sum_m\bar\sigma^{m\dot\alpha\alpha}\bar\sigma^{m\dot\beta\beta}=-2\varepsilon^{\alpha\beta}\varepsilon^{\dot\alpha\dot\beta}$ (checked exactly with the matrices of [dictionary, def:wick]), $(\Gamma(\xi,\lambda)|_V,\Gamma(\zeta,\lambda)|_V)=2\,\bar F_w$. The classical component: the bracket pairs $\lambda^*$ of $\tfrac12(\rho(F_A)\xi,\lambda^*)$ with $\lambda$ of $O_\zeta$; the $F_{A|_V}$ and $d_AA|_{W_{\mathbb R}^\perp}$ parts pair to $w^iA_{ij}w^j=0$ and to $\bar\zeta=0$; the term $-(\Gamma(\xi,\lambda)|_{W_{\mathbb R}^\perp},A^*)$ gives $(\lambda,\rho(\Gamma(\xi,\lambda)|_{W_{\mathbb R}^\perp})\zeta)$, which contains $\Gamma_{li}(\xi,\lambda)w^i=-2\varepsilon_{jkli}w^jw^i(u\lambda^k)=0$. $Q_0^\xi O_\zeta$: Theorem \ref{thm:o-invariants} (iii). The family: Proposition \ref{prop:o-laws} (b), $[s,\psi](O_\zeta)^{(1)}(\xi)=s\psi(O_\zeta)-\kappa_B\psi(B_w)$, whose $\bar F_w$ part is $-\kappa_Bx\,\bar F_w$; under $B_w\mapsto B_w+\hbar y\,P_w$, with $P_w=p_F\bar F_w+p_BB_w$ the protected combination of the two weight-$3$ singlets (the counterpart of [EJSS05, (99)], [targets, T5 (d)]), the protected direction is kept and $c_K$ changes by $-\kappa_Byp_F$.
\end{proof}

\begin{proposition}[Graphs of $c_K$]
\label{prop:o-cKgraphs}
$c_K$ is the coefficient of $\bar F_w$ in the $\varepsilon\to0$ limit of the graphs with two external legs $\bar\lambda_{\dot\alpha i}w^i$ in classes (A)--(D) of Theorem \ref{thm:o-classes} with $O_\zeta$ on the cycle, plus bridge pairs, which tend to $0$. Classes (T), (E1), (E2) contribute nothing to $\bar F_w$: $\delta_\xi O_\zeta$ is trilinear and has no one-loop graph with two external $\bar\lambda$; $\partial_\delta I^{CT}_{\mathfrak{psl}(4|4)}[\varepsilon](\xi)$ of $O_\zeta$ vanishes (degree $-1$, weight $3$, singlet: no field content); $\{I^{CT}_{\mathfrak{psl}(4|4)}[\varepsilon](\xi)+F^{(1)}(\xi),O_\zeta\}$ has outputs $(\lambda,\lambda)$ and $(A,\partial A)$. Each graph has superficial degree $0$; its value is a series in integer powers of $\varepsilon$ after the scaling $t\mapsto\varepsilon t$ of the Schwinger times (special edge in $[0,1]$, other edges in $[1,\infty)$), and the limit is the coefficient of $\varepsilon^0$ at zero external momentum. All graphs converge at $T=\infty$ (Theorem \ref{thm:o-ir} (i)).
\end{proposition}

\begin{proof}
Two external $\bar\lambda$ with both legs of $O_\zeta$ internal force $O_\zeta$ onto the cycle (a vertex with both legs joined to a connected graph lies on its cycle). Degree: output weight $3$, operator weight $\tfrac52$, $S^{(1)}(\xi)$ of weight $4$ with $\xi$ of weight $-\tfrac12$ [ch07, prop:c7-select]: superficial degree $\tfrac52+\tfrac12-3=0$. Integer powers: as in [ch04, lem:c4-special].
\end{proof}

The engine computations, in this order (details, sectors, checks and resources in the engine plan). Every reduction of an output uses full $s_{\mathfrak{psl}(4|4)}$-exact cochains with their arity partners (Remark \ref{rem:o-exact}).
\begin{itemize}
\item E0. Conventions in code: the twisted coefficient $-\tfrac12$ of Lemma \ref{lem:o-ev} (done for $\mathfrak{su}(2)$); the décalage signs of hypothesis (H1) for products and for the $C^\bullet$-linear extension of $L_1$; the tail \eqref{eq:o-Kir} (done); $\kappa_B$; the identity $\tfrac12(\Gamma(\xi,\lambda)|_V,\Gamma(\zeta,\lambda)|_V)=\bar F_w$; $B$ and $G_B$ in the realisation (derivation property done); output bases from the selection rules for every target.
\item E1. $c_K$ (Propositions \ref{prop:o-konishi}, \ref{prop:o-cKgraphs}); then $c_B'$ from the three-leg graphs.
\item E2. The family of $c_K$: under $\psi(B_w)=x\,\bar F_w$, $c_K\mapsto c_K-\kappa_Bx$; under field redefinitions, the coupling direction and $G_B$, $c_K$ is unchanged ($G_B$ shifts the $B_w$ coefficient by $-a\kappa_B$); reported as a family.
\item E3. Action side along $\xi$: $I^{CT}_{\mathfrak{psl}(4|4)}[\varepsilon](\xi)$ with two legs (no $\varepsilon^{-1}$ by weight; test the prediction $\log\varepsilon$ times the unit times $\{S^{(1)}(\xi),2\int(A,A^*)-4\int(\lambda^*,\lambda)-\int(c^*,c)\}$); the arity-one obstruction along $\xi$ in the sectors $(\lambda,A|_V)$, $(\lambda,A|_{W_{\mathbb R}^\perp})$, $(c,\lambda^*)$, including $\{S^{(1)}(\xi),F_0\}$; $F^{(1)}(\xi)$ with its family, which contains $a\,S^{(1)}(\xi)=a\,G_B^{(1)}(\xi)$; E3b: Proposition \ref{prop:o-arity0} and hypothesis (H4).
\item E4. The traceless bilinear in the plane of $w$: $Q_1=0$ (Proposition \ref{prop:o-planevanish}); a general traceless bilinear along $\xi\in\Sigma$: one classical-type coefficient, reported as a family under $G_B$ and the field redefinitions.
\item E5. The Konishi primary along $\xi\in\Sigma$: the three numbers of [ch07, cor:c7-Q], with their families.
\item E6. Letters along $\xi$, arities $0$ and $1$.
\item E7. Pair entries [targets, T4, T6]; the arity-two $(\xi,\xi)$ graphs; the two representatives of $b$.
\item E8. $S_1$: \eqref{eq:o-SQ} on the bilinears; $U_\eta$ on $B_w$, $\bar F_w$; the commutant of $\xi$.
\item E9. The arity-two $(\xi,\eta)$ component and its dilatation part $\tfrac12(\xi,\eta)$ times $D_1$ (Part D).
\item E10. $K$ at finite $T$: the tree $Y$ of Theorem \ref{thm:o-ir} (iii) in the sector $(A,A^*)$ and its arity-two analogue for the $P$-edge pairs (a tree with one constant and one linear $c$-stub, proportional to $v_1$); $K_1$ by Definition \ref{def:o-K} on the bilinears at two or more values of $T$, testing Proposition \ref{prop:o-Kclaim}, compared with \eqref{eq:o-KD}; the cross-checks of Proposition \ref{prop:o-Kconsistent}.
\item E11. $P_1=M_1=R_1=0$: internal insertions and arity-one counterterms along the covariant generators.
\item E12. The rank-$(1,1)$ element $x=\xi+\bar\xi'$, $\bar\xi'\in S_-\otimes W^*$ of rank one with $\bar\xi'_{\dot\alpha i}w^i=0$: the components along $\bar\xi'$ and along $(\xi,\bar\xi')$, and the class statement $[Q_1^x]=0$ on $H(Q_0^x)$ [targets, T12].
\end{itemize}

\begin{proposition}[The traceless bilinear in the plane of $w$]
\label{prop:o-planevanish}
Let $w_1,w_2\in W_{\mathbb R}^\perp$ satisfy $w_{a\,ij}w^j=0$. Under the hypotheses of [ch07, prop:c7-select], the component along $\xi=u\otimes w$ of $O=\int\omega\,(\iota_{w_1}A,\iota_{w_2}A)$, two-component $\tfrac1{16}\int\omega\,(w_1^{ij}A_{ij},w_2^{kl}A_{kl})$, vanishes at every order at which the hypotheses hold, in every equivariant scheme.
\end{proposition}

\begin{proof}
$(w_1,w_2)=\tfrac14w_1^{ij}w_{2ij}=0$, since $w_1^{ij}$ is a combination of $w^{[i}m^{j]}$ and $w^iw_{2ij}=0$. By [ch07, cor:c7-Q] the component is a multiple of $\int\omega\,[(\iota_{w_1}A,(\rho(w_2)\xi,\lambda))+(\iota_{w_2}A,(\rho(w_1)\xi,\lambda))]$, and $(\lambda,\rho(w_a)\xi)=(w_a,\Gamma(\xi,\lambda))=\tfrac14w_a^{kl}\Gamma_{kl}(\xi,\lambda)=-\tfrac12w_a^{kl}\varepsilon_{jikl}w^j(u\lambda^i)=0$ [ESW, (28)], (G2). Under $G_B$ the component changes by $-a$ times the classical component, which is the same functional and vanishes as well.
\end{proof}
